\documentclass[11pt,a4paper]{article}
\usepackage[T1]{fontenc}
\usepackage[utf8]{inputenc}
\usepackage{lmodern,microtype}
\usepackage[margin=.85in]{geometry}
\usepackage{amsmath,amssymb,amsthm,mathtools,aliascnt}
\usepackage{booktabs,tabularx,array,enumitem,needspace,float}
\usepackage[section]{placeins}
\usepackage{xcolor,tikz,mdframed,graphicx}
\usetikzlibrary{arrows.meta,positioning,calc,decorations.pathreplacing}
\usepackage[hidelinks]{hyperref}
\usepackage[nameinlink,noabbrev]{cleveref}
\usepackage{fancyhdr}
\hypersetup{pdftitle={Sync HotStuff RR, Part II: Sequencer Rotation},pdfsubject={Regional rotation with Espresso ordering and explicit L1 multisignature control}}
\setlength{\headheight}{14pt}
\pagestyle{fancy}\fancyhf{}
\lhead{\small Sync HotStuff RR, Part II}\rhead{\small Sequencer Rotation}
\cfoot{\thepage}
\setlength{\parindent}{0pt}\setlength{\parskip}{4pt}
\setlength{\emergencystretch}{2em}
\renewcommand{\topfraction}{.92}
\renewcommand{\textfraction}{.06}
\renewcommand{\floatpagefraction}{.85}
\makeatletter
\setlength{\@fptop}{0pt}\setlength{\@fpbot}{0pt plus 1fil}
\makeatother
\setlist{nosep,leftmargin=*}
\newcolumntype{Y}{>{\raggedright\arraybackslash}X}
\definecolor{rrRed}{HTML}{A51C30}
\definecolor{rrBlue}{HTML}{2563EB}
\definecolor{rrTeal}{HTML}{0F766E}
\newcommand{\RR}[1]{{\color{rrRed}#1}}
\newcommand{\QC}{\mathcal C}
\newcommand{\FC}{\mathcal F}
\newcommand{\Hash}{\mathsf H}
\newcommand{\rank}{\operatorname{rank}}
\newcommand{\TB}{\mathit{TB}}
\newcommand{\tagg}[1]{\ensuremath{\mathsf{#1}}}
\newcommand{\defterm}[1]{\textcolor{rrBlue}{\textit{#1}}}
\newtheorem{definition}{Definition}[section]
\newaliascnt{assumption}{definition}
\newtheorem{assumption}[assumption]{Assumption}
\aliascntresetthe{assumption}
\crefname{assumption}{assumption}{assumptions}
\Crefname{assumption}{Assumption}{Assumptions}
\newaliascnt{proposition}{definition}
\newtheorem{proposition}[proposition]{Proposition}
\aliascntresetthe{proposition}
\crefname{proposition}{proposition}{propositions}
\Crefname{proposition}{Proposition}{Propositions}
\newcounter{protocol}
% Protocol boxes float to the top of a page, so long boxes leave no gaps.
\newenvironment{protocol}[3][htbp]{\begin{figure}[#1]\refstepcounter{protocol}\label{#2}
\begin{mdframed}[linewidth=.6pt,skipabove=0pt,skipbelow=0pt,
innertopmargin=7pt,innerbottommargin=7pt]\small
\textbf{Protocol \theprotocol. #3}\par\smallskip}{\end{mdframed}\end{figure}}
\crefname{protocol}{Protocol}{Protocols}
\Crefname{protocol}{Protocol}{Protocols}
\title{Sync HotStuff RR, Part II: Sequencer Rotation}
\author{}\date{19 October 2026}
\begin{document}
\maketitle
\begin{abstract}
Part~I runs one synchronous BFT engine inside one region. This document
lets several regions take turns producing one chain. Each regional turn
is a \emph{session} with a fixed DATA quota.

The design keeps two chains apart. Regions produce a fast
\emph{production branch} and commit it under their own BFT rules.
Every verifier reads
the same finalized L1 controls and Espresso history, and so accepts the
same continuous prefix of DATA, the \emph{global prefix}~$G$. $G$ only
grows.

Handoffs come in two kinds. In a \emph{regular handoff}, once a session
has a complete regional receipt for its last DATA, the next region starts
at once; it does not wait for $G$. The next region is fixed in advance by
an immutable region map. A \emph{Force} is the only other kind. A
multisignature committee announces it on L1 as an atomic branch update
that takes effect at a future Espresso time and carries the complete
region map for the new branch. A Force closes the current production
branch and restarts from $G$. It never removes DATA that is already in
$G$, and missing output alone never causes one.

A Publish Layer sends complete regional results to Espresso. On each
node, a Session Manager starts and stops one program per session. The
document states which guarantees follow from the model and which
require additional regional service, publication service or committee
behavior.
\end{abstract}

\section{Overview}
\label{sec:overview}
Part~I~\cite{part1} gives one region a synchronous BFT engine.
This document connects several regional engines into one chain. Regions
take turns. Each turn is a \defterm{session}: one region produces a fixed
number of DATA blocks, its \emph{quota}, and then hands off to the next
region.

Two goals pull in different directions. Handoffs should be fast: the next
region should not wait for slow global confirmation. But the chain that
users rely on must never roll back, even when operators step in to
replace a failing region. The design meets both goals by keeping two
chains apart.

\Needspace{30\baselineskip}
\subsection{The core idea}
\label{sec:core-idea}
\begin{mdframed}[linewidth=.6pt,backgroundcolor=rrBlue!4,nobreak=true,
innertopmargin=6pt,innerbottommargin=6pt]
\begin{enumerate}[label=\textbf{\arabic*.},itemsep=4pt]
\item \textbf{Two chains: a fast one and a safe one.} Regions produce the
\defterm{production branch} $R$ and commit it under their own BFT rules.
A regional \emph{receipt}, made of FINAL signatures, makes a block final
inside its region. $R$ may run ahead. The \defterm{global prefix} $G$ is
the part that every verifier has accepted from public history. $G$ only
grows. Users and later stages rely on $G$.
\item \textbf{Regular handoff: start at once, successor fixed in
advance.} As soon as a session's last DATA, its \emph{terminal DATA}, has
a regional receipt, the next region starts after that block. It does not
wait for the block to enter $G$. The next region is computed from the
branch's immutable region map, so every node picks the same successor,
no matter when the receipt arrives.
\item \textbf{Force: explicit, scheduled, and restarted from $G$.} The
only other way to change regions is a \defterm{Force}. A multisignature
committee signs an atomic branch update on L1. It names a future Espresso
time, the branch version to close and a complete new region map. When
Espresso time reaches that point, every verifier closes the branch,
installs the new map and restarts production from the tip of $G$ at that
moment. The update does not say which DATA to keep; replay keeps exactly
$G$. DATA of the closed branch beyond $G$ is dropped. DATA in $G$ is never
removed. Silence, alarms and local timeouts never cause a Force.
\item \textbf{One public rule computes $G$.} Every verifier runs the same
deterministic replay over finalized Espresso blocks and finalized L1
updates. It appends a DATA block only if the block extends $G$, comes
from the session that owns the next height (the \emph{export session}),
and has a valid regional receipt. The same input gives the same $G$
everywhere.
\item \textbf{One program per session, started directly.} On each node, a
Session Manager runs a separate program for every session, with its own
votes, locks and timers. A new program starts from a fixed
\emph{anchor}, the block it continues after. Its members relay a
ReadyBundle so that all of them enter within one regional delay bound,
and the first leader proposes at once. When the Manager verifies a
Force, it stops every program of the closed branch.
\end{enumerate}
\end{mdframed}
In one sentence: regions run ahead under regional receipts, everyone
agrees on $G$ by replaying public history, and only a scheduled L1
instruction can cut off the part beyond $G$---never $G$ itself.

\paragraph{A running example.}
We use one example throughout. The regions are A, B, C and D, in this
cyclic order.
\begin{itemize}
\item A's session produces DATA 101--110. A commits all of them inside
its region and hands off to B by a regular handoff (idea~2).
\item B continues with DATA 111--115.
\item Meanwhile, verifiers have globally accepted only DATA up to 108
(idea~1).
\end{itemize}
Now the multisignature committee orders a Force (idea~3). When the order
takes effect:
\begin{itemize}
\item everything up to 108 stays in the chain;
\item the unconfirmed DATA 109--115 of the old branch can no longer be
accepted;
\item if B is selected again, B starts a \emph{new} program that produces
a new DATA 109 after 108;
\item B's old program, which continued after 110, shuts down separately.
The two programs never share votes, locks or timers (idea~5).
\end{itemize}
\begin{center}\small
\begin{verbatim}
Before the Force
  production branch R:  A: 101 ... 108 109 110   B: 111 ... 115
  global prefix G:         101 ... 108
After the Force
  G is kept:               101 ... 108
  new B program:                       new 109, new 110, ...
  dropped:              A's old 109-110 and B's old 111-115
\end{verbatim}
\end{center}
\Cref{fig:concurrent-handoffs} in \cref{sec:xprotocols} draws the same
example with all protocol details. The rest of the document explains when
a handoff happens, who chooses the next region, why the new program
continues after 108, and how each node starts and stops the
corresponding programs.

\subsection{How the protocol works}
This subsection adds detail to ideas 1--3.

\paragraph{Two kinds of handoff.}
A \defterm{regular handoff} happens when a session finishes its quota.
A commits its last DATA, collects a regional commitment receipt, and
gives the result to B. B checks the data, its own selection and the
execution state, and then continues, possibly before A's last
DATA is globally accepted.

A \defterm{Force} is a handoff ordered by the multisignature committee.
The committee submits an L1 instruction that names a future Espresso
time, the target branch version and one complete new status map. All
verifiers apply it at the same place: the first Espresso block whose
timestamp reaches that time. If the target is the current branch
version, they close the old branch, install the new map and continue
from the current global prefix. If the target is not current, it is
\emph{stale}, and the whole instruction does nothing. A Force happens
only through such an instruction. Silence of a region, a monitoring
alarm or a local timeout does not cause one. The committee may use slow
delivery as a reason to act, but that judgment belongs to its own
policy.

\paragraph{Three stages of a block.}
A DATA block passes through three separate stages.
\begin{enumerate}
\item \emph{Regional commitment.} The producing region commits the block
under its BFT rules and collects enough FINAL signatures.
\item \emph{Publication.} The Publish Layer sends the block and its
evidence to Espresso, which fixes its position in a finalized log.
\item \emph{Global acceptance.} Verifiers check the regional evidence,
the L1 controls and the rotation rules. If the block extends the
current global prefix, they append it.
\end{enumerate}
Only the third stage adds a block to the global prefix $G$.
We write $G_q$ for its value at Espresso height $q$, and ``$G=108$'' for
``the global prefix ends at DATA 108''.

The production branch $R$ is what regions have produced so
far. It can run ahead of $G$: in the example, $R$ reaches 115 while
$G$ is still 108. The \defterm{export session} is the session that owns
the next DATA height after $G$. In the example it is still A's
session, because 109 is A's block, even though B is the region that is
currently producing. A Force closes the part of $R$ beyond $G$; it
never touches $G$.

A fourth stage lies outside this document: rollup DA publication. It
has its own accepted prefix. That prefix may lag behind $G$, but it may
not skip any DATA in $G$. The Publish Layer of this document is a
different service: it only delivers regional results to Espresso.

A regional receipt stays valid evidence of regional commitment, even if
the branch that contains the block is closed later. Global acceptance is
stronger: no Force ever removes DATA from $G$. All these guarantees are
about the input: which DATA, in which order, from which session. Correct
execution and native state continuity are checked by the application and
publication interfaces, not here.

\begin{table}[htbp]\centering\small
\caption{Participants and responsibilities.}\label{tab:roles}
\begin{tabularx}{\linewidth}{@{}>{\raggedright\arraybackslash}p{.19\linewidth}Y@{}}\toprule
Role & Responsibility\\\midrule
L1 & Publishes configuration and recovery records and the ordered control log.\\
Regional replica & Runs the DATA/DRAIN engine of its sessions and
serves complete commitment evidence.\\
Multisignature service & Announces future atomic branch updates on L1.
Each update supplies the complete map for the new branch and causes a
Force if its target is current. It cannot choose a retained prefix or
declare data global.\\
Publish Layer & Delivers complete regional results to Espresso. The
leader, replicas and external verifiers may act as senders.\\
Espresso & Supplies a finalized ordered log, authenticated timestamps,
complete namespace contents, and bounded finalization and visibility.\\
Observer / certifier & Reads Espresso and L1, runs the global finality
rule, and so derives $G$, the next session and every Force boundary.\\
\bottomrule
\end{tabularx}
\end{table}

\subsection{The four view indices}
A replica's position in the protocol is its \emph{view}. The full view
has four indices, from coarse to fine:
\[
 (\beta,r,g,t).
\]
\begin{itemize}
\item $\beta$ is the \defterm{branch version}. Each effective Force
increases it by one. It is not an Espresso epoch or a time interval; it
only separates the old production branch from the new one.
\item $r$ is the \defterm{turn index}. Together with $\beta$, it identifies
the session whose authority fixes the voting snapshot and DATA quota.
It is not a physical region identifier.
\item $g$ is the \defterm{generation}: how many times the session has
run regional recovery.
\item $t$ is the \defterm{term}: which leader turn within the
generation.
\end{itemize}
A \defterm{session} is the pair $(\beta,r)$: one region's production
turn on one branch version. Its voting snapshot and DATA quota stay
fixed through all of its recoveries. Every node and certifier derives
these numbers by replaying the same authenticated history, so no other
session identifier is needed.

\begin{table}[htbp]\centering\small
\caption{How the view changes. Each trigger is defined in the cited
section.}\label{tab:views}
\begin{tabularx}{\linewidth}{@{}>{\raggedright\arraybackslash}p{.19\linewidth}>{\raggedright\arraybackslash}p{.25\linewidth}Y@{}}\toprule
Change & New view & Trigger and effect on timers\\\midrule
Leader rotation & $(\beta,r,g,t+1)$ & The last block of term $t$ gets a
QC (its terminal QC) while the session still has DATA quota
(\cref{sec:quota}). Running commit timers continue.\\
Regional recovery & $(\beta,r,g+1,0)$ & A quit proof freezes generation
$g$. After a $\Delta_i$ wait, and once the recovery authority is
finalized, replicas enter $g+1$ (\cref{sec:quit}). The quit cancels
generation-$g$ timers.\\
Regular handoff & $(\beta,r+1,0,0)$ & A valid receipt for the session's
terminal DATA fixes the successor; each member enters once it is Ready
(\cref{sec:witness,sec:ready}). The old engine stops through
Completion.\\
Force & $(\beta+1,r',0,0)$ & An activated branch update whose target is
$\beta$ (Part~IV~\cite{part4}). Let $r_c$ be the export session's turn index:
$r'=r_c$ if that session has no accepted DATA, otherwise $r_c+1$; an
unassigned boundary keeps its target index. Verifying the Force
terminates old-version programs and cancels their timers.\\
\bottomrule\end{tabularx}
\end{table}

For example, $(7,20,0,0)\to(7,20,0,1)$ is a leader rotation inside A,
$(7,20,0,1)\to(7,20,1,0)$ is a recovery inside A, and
$(7,20,1,0)\to(7,21,0,0)$ is a regular handoff to B. None of these
changes $\beta$.

A Force can reuse a turn index. Suppose export session $(7,20)$ has
already added DATA to $G$. A Force then opens $(8,21,0,0)$, even if
a regular successor already used $(7,21)$. If $(7,20)$ had added no
DATA to $G$, the Force would open $(8,20,0,0)$ instead. In both cases
the new pair differs from every old pair, because $\beta$ differs. So
votes, locks and evidence stay separate. Where the new session starts is
fixed by the DATA anchor it keeps, not by its turn index.

An L1 control that is finalized but not yet effective changes nothing:
production and global acceptance go on as before.
, BLAME and Quit remain in place. Inside a region,
the only changes from Part~I are those listed in \cref{sec:protocols}.


\section{Model and assumptions}
\label{sec:model}
This section fixes what the protocol relies on: what L1 records, how
results reach Espresso, what faults are allowed, and which timing and
service bounds hold.

\subsection{L1 authorities and atomic branch updates}
\label{sec:l1}
\label{sec:runtime-controls}
L1 plays two roles. First, it stores the configuration and recovery
records used to derive immutable session \emph{authorities}. Second,
it stores the \emph{control log}: an ordered list of
atomic branch updates signed by the committee. A branch update is the
only way to change the region map, and the only way to force a
handoff.

Nodes authenticate finalized L1 headers, authority source records and a
complete, ordered control history. L1 timestamps only label L1 blocks. They are
neither regional commit times nor Espresso activation times.

\begin{definition}[\defterm{Atomic branch update}]
\label{def:control-update}
An atomic branch update is
\[
 U=(\mathit{chain},c,\theta_a,\beta_{\rm target},S_{\rm new})
\]
together with a signature set $\Sigma_U$ on its hash. Here $c$ is its
sequence number, $\theta_a$ its Espresso activation threshold,
$\beta_{\rm target}$ the branch version it asks to close, and
\[
 S_{\rm new}:\mathcal L_R\longrightarrow
 \{\tagg{ON},\tagg{OFF}\}
\]
is one explicit, complete status map; it may set every region OFF, but
the map field itself may not be omitted. The update names no retained
DATA height. There are no status-only updates and no separate Force tag.
\end{definition}

\begin{definition}[\defterm{Valid atomic branch update}]
\label{def:control-valid}
An update $U$ in finalized L1 history is valid when:
\begin{enumerate}[label=(V\arabic*)]
\item $\mathit{chain}$ is this chain;
\item $\Sigma_U$ satisfies the control authority $\mathcal K$ fixed by
the genesis checkpoint;
\item $c$ is larger than the sequence number of every earlier valid
update in L1 order;
\item $\beta_{\rm target}\in\mathbb N$ and $S_{\rm new}$ assigns exactly
one flag to every region of $\mathcal L_R$; and
\item $\theta_a\ge\theta_a(U')+T_{\rm update}$, where $U'$ is the
previous valid update, or the last update of the genesis prefix.
\end{enumerate}
\end{definition}

The \defterm{control log} $\Lambda=\langle U_1,U_2,\ldots\rangle$ is the
sequence of valid updates in L1 order. A verifier processes every entry
exactly once, including stale entries, and keeps the last processed
sequence and control-prefix commitment as its control cursor $\kappa$.

\begin{table}[htbp]\centering\small
\caption{L1 authority and control updates.}\label{tab:l1}
\begin{tabularx}{\linewidth}{@{}p{.27\linewidth}Y@{}}\toprule
Field & Meaning\\\midrule
$\tagg{AUTH}_{\beta,r,0}$ & Initial session authority: fixed configuration,
leader-rule inputs and DATA quota $\TB_r$.\\
$\tagg{AUTH}_{\beta,r,g}$, $g>0$ & Initial authority plus the finalized
recovery record for $(\beta,r,g)$; the snapshot and quota stay fixed.\\
Branch update & One L1 record
$(\mathit{chain},c,\theta_a,\beta_{\rm target},S_{\rm new})$
(\cref{def:control-update}).\\
Control sequence $c$ & Unique increasing index binding the chain and the
update.\\
Activation time $\theta_a$ & Espresso timestamp threshold, announced
before its first crossing can become final. One per atomic update;
strictly increasing along the log, at least
$T_{\rm update}$ apart.\\
Target $\beta_{\rm target}$ & The branch version that the update asks to
close. A stale target makes the entire update a no-op.\\
Status map $S_{\rm new}$ & A complete ON/OFF map for every region. If
the update executes, this is the immutable map of the new branch
version.\\\bottomrule
\end{tabularx}
\end{table}

\paragraph{Authority contexts and leaders.}
\label{sec:authority-context}
$\tagg{AUTH}_{\beta,r,g}$ is the immutable authority for generation $g$
of session $(\beta,r)$, bound to its authenticated selection context.
Its initial configuration and leader-rule inputs are derived from
finalized L1 records and authenticated predecessor or Force history.
These inputs are fixed before target-region selection; the handoff
separately binds the selected region and exact execution anchor.
Initial authority derivation requires no new L1 transaction at a regular
handoff. Part~III~\cite{part3} supplies a staking-based construction.

$\tagg{AUTH}_{\beta,r,0}$ fixes $\mathit{cfg}_{\beta,r}$, its configuration
ID, any selection reference and seed used by the leader rule, and $\TB_r$.
The configuration contains each region's members, weights, $F_i$,
$\Delta_i$ and $K$. For $g>0$, the authority combines this initial
context with the unique finalized L1 recovery record for $(\beta,r,g)$,
bound to the same initial authority. Only leader eligibility and the
schedule may change. For each admitted term $t$ of the assigned region
$i_{\beta,r}$, the leader is
\[
 L(\beta,r,g,t)=\tagg{Leader}
 \bigl(\tagg{AUTH}_{\beta,r,g},i_{\beta,r},t\bigr).
\]
The rule uses the full view; there is no
separate L1 authority record for each term. Reusing $r$ on a new branch
does not reuse the old authority. Underlying snapshot bytes may be shared
when both contexts select them. The quota remains indexed by $r$;
controls and slashing never reweight historical certificates or reset it.

Write $\lambda(\tagg{AUTH}_{\beta,r,g})$ for the greatest L1 publication
or activation height required by its source records, including the
applicable branch-control publication and, for $g>0$, its recovery record.
This is a source-height bound, not a transaction publishing each derived
context. Missing source evidence causes a wait. Verification checks the
session, configuration and handoff bindings, not just the numeric indices.

\paragraph{Ingesting controls.}
Finalization on L1 only announces an update. The update takes effect at
its activation time. Activation times strictly increase along the log,
at least $T_{\rm update}$ apart (\cref{def:control-valid}), so entries
are processed in sequence order (\cref{def:control-effect}). Malformed,
unauthorized or duplicate updates grant nothing
(\cref{def:control-valid}). To change an instruction, the committee
submits a later update; finalized history is never edited.

An executing update atomically closes its target branch, increments the
branch version and installs $S_{\rm new}$ as the new version's immutable
status map.  If its target is stale, the entire update is a no-op: its
map does not apply and its operational intent is not carried forward.
The committee must resubmit that intent against the current version.

\paragraph{Espresso positions and when controls apply.}
Let $J_q$ be Espresso block $q$, with authenticated timestamp $\theta_q$.
Within a block, namespace-transaction and record indices give a complete
order, also when timestamps are equal. For a threshold $d$, let
\[
 q(d)=\min\{q:\theta_q\ge d\}
\]
be the first block whose timestamp reaches $d$. For example, if blocks
500, 501, 502 have timestamps 119, 122, 125, then $q(120)=501$.

A control with threshold $d$ takes effect at the start of $J_{q(d)}$,
before any EXPORT in that block is examined. (An EXPORT is a Publish
Layer record that carries regional results; see \cref{sec:delivery}.) We
call this start-of-block step the \defterm{prelude}. By
\cref{ass:single-activation}, at most one update activates in a prelude.
It executes exactly when its target equals the branch version at the
start of the prelude. Otherwise the whole update is stale and changes
nothing. There is no other path to a cross-region Force. We write ``$G$
at the start of $J_q$'' for the global prefix just before the block's
EXPORTs.

\begin{definition}[\defterm{Effect of the activated update}]
\label{def:control-effect}
Let $\beta_0$ be the branch version at the start of $J_q$. If
$\Lambda_q=\varnothing$, the prelude changes no control state. Otherwise,
by \cref{ass:single-activation}, $\Lambda_q=\{U_q\}$ for one update
\[
 U_q=(\mathit{chain},c_q,\theta_a,\beta_{\rm target},S_{\rm new}).
\]
If $\beta_{\rm target}\ne\beta_0$, $U_q$ is a stale whole-update no-op.
If $\beta_{\rm target}=\beta_0$, the prelude atomically closes $\beta_0$,
sets $\beta:=\beta_0+1$, fixes
\[
 S_{\beta}:=S_{\rm new},
\]
and executes the Force transition of \cref{alg:global-rule}. In either
case, the single update advances the processed control cursor exactly
once.
\end{definition}

Genesis fixes $S_0$.  The map $S_\beta$ is immutable for the lifetime of
branch version $\beta$; regular handoffs change $r$ but continue to use
the same map.  A region is eligible for target session $(\beta,r)$ exactly when
$S_\beta$ sets it ON and $\tagg{AUTH}_{\beta,r,0}$ gives it positive total
weight. A monitoring report, missing output or local timeout changes
nothing.  If every region is ineligible, the next session stays pending
until another atomic branch update creates a later version with a new
complete map.  No update removes accepted DATA.

\paragraph{Choosing the next region.}
$\tagg{NextLive}(i,S_\beta;\beta,r)$ picks the region for target session
$(\beta,r)$, starting after source region $i$. Eligibility is judged by $S_\beta$ and
$\tagg{AUTH}_{\beta,r,0}$, as above. Selection uses only its already fixed
configuration inputs; it does not wait for the selected region's Ready.
The selector returns the first eligible
region after $i$ in the fixed cyclic order $\mathcal L_R$, not counting
$i$ itself. If no other region is eligible, it returns $i$ when $i$ is
still eligible, and $\bot$ otherwise. A synthetic source $\bot$ starts
the search from the beginning of $\mathcal L_R$. Because $S_\beta$ is
fixed, every regular handoff in version $\beta$ obtains the same result,
no matter when its terminal receipt is delivered.

\paragraph{Genesis.}
The system starts from an initial checkpoint. It fixes the chain
identifier, the cyclic region order $\mathcal L_R$, the genesis DATA
anchor, the first replay position, the Espresso quorum trust anchor,
the control authority $\mathcal K$, the minimum update gap
$T_{\rm update}$, the control prefix, the complete initial map $S_0$,
and the initial branch version $\beta=0$. The initial
authority and state are either trusted or verified from their stated
sources; starting this replay does not make Espresso authenticate
earlier application history.

The genesis anchor carries synthetic commitment evidence and fixed,
authenticated L1 and Espresso references. At replay position
$q_{\rm init}$, a synthetic boundary with no source region and target
index 0 applies $\tagg{NextLive}(\bot,S_0;0,0)$. If it returns a region,
session $(0,0)$ opens with $q^{\rm open}_{0,0}=q_{\rm init}$; otherwise
the same boundary stays pending. A later successful atomic update closes
the pending version, installs a new complete map and tries the same target
index on the new version. In either case,
$\lambda^{\rm open}_{\beta,0}=\max\{\lambda(\text{genesis}),
\lambda(\tagg{AUTH}_{\beta,0,0})\}$, where the second term is the
authority source-height bound defined above. A later checkpoint used as a
starting point must instead include its current boundary, its reference
bounds and the complete replay state.

\subsection{Data delivery and the Publish Layer}
\label{sec:delivery}
The \defterm{Publish Layer} is a delivery service run by existing
participants. It runs no consensus of its own and does not decide which
DATA is valid. A regional leader, another replica or an external
verifier can submit. The rotation rules rely only on the service
guarantee in \cref{ass:services}(i); Part~IV~\cite{part4}
describes one implementation.

A \defterm{complete regional result} is a DATA block together with its
direct or covering FINAL receipt and the ancestry, boundaries and source
evidence needed to validate it. A proposal QC, a block hash, or a receipt
without its payload is not complete. Correct replicas offer complete
results to the Publish Layer as soon as they have them. The service may
batch results and builds EXPORT records as specified in
\cref{sec:exports}.

Regions may also send results directly to other regions. This lets a
regular successor become Ready before the terminal result is finalized
in Espresso. Without direct delivery, the successor reads the result
from Espresso after finalization, which takes longer. There is no
separate cross-region transport service.

Espresso carries finalized records, including DATA and commitment
evidence, to online readers. Readers still check the application and
rotation rules. L1 supplies configuration, recovery and control proofs. An
execution-state snapshot may be transferred directly or rebuilt from
authenticated input history. Rebuilding state is part of Ready,
not of Espresso inclusion.

\subsection{Stake and faults}
\label{sec:stake}
Member $v$ of region $i$ has voting weight $\omega_{i,v}>0$. Let
$W_i=\sum_{v\in V_i}\omega_{i,v}$. The stake configuration fixes a bound
$F_i<W_i/2$, and Byzantine members of region $i$ hold total weight at
most $F_i$. Every regional certificate needs matching signatures from
distinct members with total weight greater than $F_i$. So every
certificate has a correct signer, and the correct members alone can form
one. This replaces the $f+1$ of $n=2f+1$ in Part~I.

Membership and weights stay fixed for a whole session, including its
recoveries. Slashing and later atomic branch updates do not change them. A later
session checks old certificates against the old snapshot. The bound
$F_i$ counts every key of the snapshot that is compromised while its
certificates are still accepted. Compromise of old keys after that is
outside this model.

Correct replicas keep their protocol state and committed blocks.
Signatures cannot be forged, and hashes are binding.

\subsection{Assumptions}
\label{sec:assumptions}
In short, the assumptions below say: finalized L1 and Espresso histories
are common and never change; messages inside a region arrive within
$\Delta_i$; the committee is honest; at most one update activates per
Espresso block; and the Publish Layer, L1 and Espresso deliver within
fixed bounds. These bounds never trigger a Force.
\begin{assumption}[Finalized L1 records]
\label{ass:authority}
At every finalized L1 height, the configuration and recovery records
and ordered control history through that height are common and
authenticated. Finalized records never change. The same authenticated
session history and source records determine the same authority context.
A verifier obtains a complete control
prefix, not an arbitrary subset of valid events.
\end{assumption}

\begin{assumption}[Authenticated synchronous communication and global clock]
\label{ass:network}
\leavevmode\par
\begin{itemize}
\item Messages use authenticated communication channels.
\item Within region $i$, the actual delay of a correct-to-correct
message is $\delta\le\Delta_i$; $\delta$ may differ across messages.
\item All nodes share a global clock.
\end{itemize}
\end{assumption}
The delay bound covers delivery and processing of complete validation
dependencies. Each replica records its own entry and voting events on
the global clock.

\begin{assumption}[Honest multisignature committee]
\label{ass:committee}
The committee acts as one honest, coordinated control service. It
submits only authorized, immutable, future-dated atomic branch updates, with
activation times that satisfy the advance-announcement condition below. It keeps no private protected prefix and cannot select
an older retained anchor. Honesty here imposes no response deadline; the
minimum gap $T_{\rm update}$ between updates is a validity rule
(\cref{def:control-valid}), not part of honesty.  The reference policy
normally keeps one unactivated update in flight and resubmits any stale
operational intent against the current branch version.
\end{assumption}

\begin{assumption}[At most one activation per Espresso block]
\label{ass:single-activation}
For every Espresso height $q$, at most one valid atomic branch update
$U$ satisfies $q(\theta_a(U))=q$. Equivalently, the activated-update
input $\Lambda_q$ is either empty or a singleton. This is a model
assumption; it is not derived from the gap $T_{\rm update}$.
\end{assumption}

\Needspace{8\baselineskip}
\begin{assumption}[Espresso history and time]
\label{ass:espresso}
Finalized Espresso history is consistent, persistent and publicly
verifiable, including the complete namespace content of each block, so a
verifier can also confirm that a block holds no record for this chain.
Each block carries a non-decreasing timestamp that runs at most
$\epsilon_E\ge0$ ahead of real time: if $J_q$ first becomes final at real
time $t_E(q)$, then $\theta_q\le t_E(q)+\epsilon_E$.
\end{assumption}
Long-term retrieval is a separate service, for example archival readers;
a finality certificate alone does not provide it. A returning node
catches up, or starts from an authorized checkpoint, before it confirms
new output; catching up has no fixed time bound.

\begin{assumption}[Bounded external services]
\label{ass:services}
There are fixed bounds $T_{\rm pub},T_{\rm L1},T_E>0$, valid under the
declared input-size, admission and load limits, such that:
\begin{enumerate}[label=(\roman*)]
\item \emph{Publish Layer.} If a correct participant has a complete
regional result available at real time $s$, an Espresso-admissible record
carrying that result, with its DATA and commitment evidence intact, is
submitted to Espresso by $s+T_{\rm pub}$. Submission does not rely on the
original leader alone: a working submission path and the required fee
resources exist.
\item \emph{L1.} A valid control operation actually submitted at real
time $s$ is finalized by $s+T_{\rm L1}$, and by then every correct online
participant can retrieve and verify it.
\item \emph{Espresso.} An admissible complete record submitted at real
time $s_E$ is in a finalized block by $s_E+T_E$, and by then every correct
participant that stays online and keeps up with the history can retrieve
the complete record and verify its inclusion and finality. Every finalized
block, with its complete relevant namespace, becomes retrievable and
verifiable to those participants within $T_E$ of finalization, also when
it holds no record for this chain.
\end{enumerate}
\end{assumption}
Participants are regional replicas, external observers and certifiers.
Each bound covers the whole path: inclusion, finalization, dissemination,
retrieval and validation. Each clock starts at an actual event: in (i),
when the complete result is available, not when a leader proposes the
block or one replica commits it; in (ii) and (iii), at actual submission,
which excludes any detection or decision time before it. ``Submitted''
means the complete record reaches the submission interface under its
admission rules; a local send attempt is not enough. A result already
recorded in Espresso satisfies (i). Forming the FINAL receipt is the
regional protocol's job.

These are delivery bounds only. Espresso delivers full payloads and
evidence, but does not promise that an EXPORT passes this chain's
acceptance rules: a stale base, a closed branch or a missing predecessor
can still stop it from advancing $G$, and retries needed for acceptance
are a separate task. The bounds never trigger a Force, and $T_E$ does not
replace the clock-lead bound of \cref{ass:espresso} or the
advance-announcement rule below.

For a newly available complete result, (i) and (iii) give
\[
 s_{\rm submit}\le s+T_{\rm pub},\qquad
 s_{\rm visible}\le s+T_{\rm pub}+T_E,
\]
where $s_{\rm visible}$ is the time by which the result and its Espresso
evidence are available to the readers in (iii). This bounds neither when
the regional result first forms, nor when a new region restores its
execution state.

\paragraph{Advance announcement.}
A control must be visible to everyone before it can take effect.
For a control submitted at real time $s$, the committee chooses
\[
 \theta_a>s+T_{\rm L1}+\epsilon_E.
\]
Then the control is available to correct participants before the first
block crossing $\theta_a$ can become final. This is a model condition,
not a conclusion from average L1 or Espresso block times. The L1 bound
already includes dissemination, so no second announcement delay is
added. $T_{\rm L1}$ starts at actual submission; detection, deliberation
and signature collection before submission are not covered by any
bound. Offline or unsynchronized participants catch up before they
confirm new output. Missing control history pauses verification.

\paragraph{Verification interfaces.}
We assume that regional receipt validity and the application-validity
checks below are sound under the regional model. A deployed certifier
must authenticate both source histories and bind the control version it
used. A signature or a TEE alone does not show that an old control
snapshot is complete. The publication interface additionally requires
sound execution and DA correspondence checks.

The bounded preparation assumption for a newly selected region is
stated with the Ready rules, in \cref{ass:ready-service}. The conditions
needed for progress are collected in \cref{sec:progress}.

\section{Regional engine: definitions}
\label{sec:regional}
This section defines the objects that one regional engine uses: blocks
and their positions, the session quota, certificates and receipts, and
the rules for admitting views. The engine is Part~I's engine with
sessions, a quota, DRAIN blocks and FINAL receipts added.

Three objects matter most for rotation: the \emph{anchor} that a session
continues after; the \emph{terminal DATA}, the session's last DATA block,
fixed by its quota; and the \emph{regional receipt} for a DATA block,
which later sessions and verifiers use as proof of regional commitment.
\Cref{tab:notation} lists the main symbols. Anchors and entry times
carry the same four indices as the view.
\begin{table}[htbp]\centering\small
\caption{Notation.}\label{tab:notation}
\begin{tabularx}{\linewidth}{@{}p{.32\linewidth}Y@{}}\toprule
Symbol & Meaning\\\midrule
$(\beta,r,g,t)$; $j$ & Branch version, turn index, generation, term; fresh position.\\
$(\beta,r)$ & Session: one region\textquotesingle s turn on one branch version.\\
$i_{\beta,r}$; $V_i$ & Session region; its regional members.\\
$\omega_{i,v},W_i,F_i$ & Member weight, total weight, and Byzantine bound $F_i<W_i/2$.\\
$\Delta_i$; $K$ & Regional synchronous delay bound; maximum DATA per term.\\
$\tagg{AUTH}_{\beta,r,g}$; $\mathit{cfg}_{\beta,r}$ & Immutable authority; fixed session snapshot.\\
$A_{\beta,r,g,t}$ & Anchor: the block just before the first fresh position of that view.\\
$h(B)$; $k_{\beta,r}(B)$ & DATA height; physical position in the session, counting DRAIN.\\
$\TB_r,H_{\beta,r},T_{\beta,r}$ & Session DATA quota, last DATA height, and terminal DATA.\\
$m_{\beta,r,g,t}$; $L(\beta,r,g,t)$ & Number of fresh positions in the term; designated leader.\\
$\QC_{\beta,r,g,t}(B)$; ANCHOR QC & Proposal certificate; re-certification of the opening.\\
$\FC_{\beta,r}(B)$ & Regional FINAL receipt, or its same-session coverage.\\
$Q_{\beta,r}^0,Q_v,Q^*$ & Inherited commitment evidence; entry lock; selected justification.\\
$e_v(\beta,r,g,t),\tau_v(P),a_v$ & Real-time entry, vote and quit events.\\
$\rho,\lambda(B)$ & Finalized L1 reference and its L1 height.\\
$\xi(B)$ & Signed Espresso header reference.\\
$J_q,\theta_q$ & Espresso block and authenticated timestamp.\\
$S_\beta$; $\mathcal L_R$ & Immutable status map of branch version $\beta$; cyclic region order.\\
$q_{\beta,r}^{\rm open}$ & The session's lower bound for Espresso references; not a successor-selection time.\\
$\lambda_{\beta,r}^{\rm open}$ & The session's lower bound for L1 references.\\
$G$, $G_q$; $R$ & Global DATA prefix, and its value at Espresso height $q$; production branch.\\
$T_{\rm L1},\epsilon_E$ & L1 announcement-service bound; maximum Espresso clock lead.\\
$T_{\rm pub},T_E$ & Complete-result submission bound; Espresso finalization and visibility bound.\\
$T_{\rm ready}$ & Preparation bound, measured from the public opening position.\\
$\nu=(\beta,r,i,A,H,q^{\rm open},\mathit{cause})$ & Complete descriptor of the session responsible for the next DATA.\\
$\sigma_G(B)$ & Session record $\nu$ bound to an accepted non-genesis DATA block $B\in G$.\\
$\mathsf{Tail}$; $\omega$ & Tail $\tagg{NEXT}(\nu)$ or $\tagg{PENDING}(\omega)$; boundary with no selected region.\\
$a(G,\nu)$ & Number of DATA of session $\operatorname{sid}(\nu)$ already accepted into $G$.\\
$U=(\mathit{chain},c,\theta_a,\beta_{\rm target},S_{\rm new})$; $\mathcal K$ & Atomic branch update; control authority that signs it.\\
$S_\beta$; $\Lambda$; $\Lambda_q$; $T_{\rm update}$ & Immutable map of branch $\beta$; control log; empty-or-singleton update activated at $J_q$; minimum activation gap.\\
$\mathcal X_q$; $\kappa_q$ & Maintained replay state; processed control cursor.\\
$\mathit{st}(B)$ & Executable state after block $B$.\\
$E$; Ready & Regional entry proof; the Ready condition of \cref{def:ready}.\\
\bottomrule\end{tabularx}
\end{table}

\subsection{Blocks, positions, and L1 references}
\label{sec:blocks}
The \defterm{anchor} $A_{\beta,r,g,t}$ is the block just before the first
fresh position of view $(\beta,r,g,t)$, that is, before the first new
block the view proposes (\cref{sec:quota}). In short: the anchor is the
block the view continues from. We always write all four indices.
\begin{itemize}
\item $A_{\beta,r,0,0}$ is the DATA block inherited from the preceding
session; the first session inherits genesis. It keeps the regional
commitment evidence of the session that produced it: a direct or
covering receipt (\cref{sec:coverage}). Genesis has synthetic evidence.
\item For $g>0$, NEWVIEW chooses the \defterm{recovery anchor}
$A_{\beta,r,g,0}$ (\cref{alg:recovery}). It is certified but may be
uncommitted.
\item For $t>0$, $A_{\beta,r,g,t}$ is the terminal DATA block of term
$(\beta,r,g,t-1)$, certified by that term's terminal QC.
\end{itemize}
The anchor of a view never changes. A generation starts at
$A_{\beta,r,g,0}$, and its later terms use their own anchors. ANCHOR
voting re-certifies only $A_{\beta,r,g,0}$, once per recovery
generation $g>0$. Generation zero uses the fixed handoff directly.

In the running example, A's session starts after anchor 100, and B's
regular session starts after anchor 110. After the Force, B's new session
starts after anchor 108. The inherited anchor keeps its original
producing session, height and hash: DATA 108 produced by $(7,20)$ stays
a $(7,20)$ block even when $(8,21)$ continues from it.

Each block has two positions.
\begin{itemize}
\item The \defterm{DATA height} $h(B)$ counts the DATA blocks on the chain
from genesis to $B$.
\item The \defterm{physical position} $k_{\beta,r}(B)$ counts DATA and DRAIN
blocks within session $(\beta,r)$, with $k_{\beta,r}(A_{\beta,r,0,0})=0$.
This local count does not change the signed fields of $A_{\beta,r,0,0}$.
\end{itemize}
A new block of session $(\beta,r)$ has the form
\[
 B=(\beta,r,k,h,\mathit{type},x,\Hash(B_{\rm parent}),\rho,\xi),
 \qquad \mathit{type}\in\{\tagg{DATA},\tagg{DRAIN}\},
\]
with $k=k_{\beta,r}(B_{\rm parent})+1$. A \defterm{DATA block} carries an opaque application payload
$x$ and has $h=h(B_{\rm parent})+1$. A \defterm{DRAIN
block} carries no batch and has $h=h(B_{\rm parent})$. So DATA advances
both positions, and DRAIN advances only the physical position.

\paragraph{L1 references.}
The \defterm{reference} $\rho$ names a finalized L1 block. It holds that
block's height $\lambda(B)$ and hash, and a registry proof for the L1
control records up to that height. It tells validators which fixed
authorities and rules apply. These rules hold:
\begin{itemize}
\item along DATA ancestry, $\lambda$ never decreases;
\item the first DATA block of session $(\beta,r)$ has $\lambda\ge
\lambda^{\rm open}_{\beta,r}$, the L1 lower bound fixed by its boundary
(\cref{sec:witness}; genesis uses \cref{sec:l1});
\item a DRAIN block copies the reference of $T_{\beta,r}$.
\end{itemize}
The timestamp of the L1 block named by $\rho$ is readable from its
header, but it is not a commit time for $B$: recovery may commit $B$
much later (\cref{sec:late-receipt}). The Vote step applies no
receive-time freshness filter to $\rho$.

The L1 reference serves regional validation only. It does not choose
the time of a Force or the Espresso position used for selection.
Regional validity uses the fixed session authority, so a later status
update does not change an old receipt's voting weights.

\paragraph{Espresso references.}
The outer field $\xi(B)$ names a finalized Espresso header $J_q$ by
height and hash. Its height does not decrease along DATA ancestry, and
the first DATA of session $(\beta,r)$ has
$q\ge q_{\beta,r}^{\rm open}$. Terminal DATA uses the same rule as other
DATA. DRAIN copies the Espresso reference of $T_{\beta,r}$, just as it
copies its L1 reference. The proposal hash and signatures bind these
references and the session context. They do not change the native
mini-block encoding. Eligibility is not read from $J_q$: every session
in branch version $\beta$ uses the fixed map $S_\beta$.

Regional full validation authenticates the referenced history, directly
or from its validated cache. The public receipt check \tagg{CheckFINAL}
(\cref{sec:receipt-kernel}) establishes the intrinsic regional validity
of the reference. Global acceptance separately authenticates the
required Espresso context. It also requires the reference to name an
Espresso block strictly before the block that contains the EXPORT, so no
reference can justify itself through that same EXPORT. A historical
reference says nothing about freshness when it is received. The
reference of $T_{\beta,r}$ contributes only to the regular successor's
reference floor (\cref{sec:witness}).

\paragraph{Why the Espresso reference is needed.}
The terminal reference does not choose the next region. That choice is
already determined by the immutable map $S_\beta$, so two nodes that
receive the terminal receipt at different times compute the same
successor. Instead, the terminal reference gives the successor a common
lower bound for the Espresso headers that its first DATA may cite. It is
not proof that the terminal DATA is globally final.

For example, let the region order be A, B, C and let $S_7$ set all three
regions ON. If A's terminal DATA 110 binds
$\xi(110)=(500,\Hash(J_{500}))$, every verifier selects B under $S_7$,
whether it receives the terminal receipt before or after $J_{505}$.
Changing B to OFF requires an atomic update that closes version 7 and
creates a new version with a new map; it cannot retroactively change the
regular boundary inside version 7.

The next regular session uses the greater of two heights as its reference
floor:
\[
 q^{\rm open}_{\beta,r+1}
 =\max\{\operatorname{height}(\xi(T_{\beta,r})),
         q^{\rm open}_{\beta,r}+1\}.
\]
This advances successive floors without delaying successor selection.
The number is a lower bound, not necessarily the height of an already
available header and not the time at which the successor is selected.
The successor may be assigned and enter from the direct terminal receipt;
its first DATA must wait until it can cite a finalized header at or above
the floor. DRAIN keeps the terminal DATA's reference.

\subsection{Session quota, term positions, and DRAIN}
\label{sec:quota}
\paragraph{1. Fix the session endpoint.}
Session $(\beta,r)$ may add $\TB_r\ge1$ DATA blocks after
$A_{\beta,r,0,0}$. Its \defterm{last DATA height} is
\[
 H_{\beta,r}=h(A_{\beta,r,0,0})+\TB_r.
\]
The DATA block at this height is the \defterm{terminal DATA}
$T_{\beta,r}$. Recovery never changes $H_{\beta,r}$. The quota counts the
DATA blocks after $A_{\beta,r,0,0}$ in a proposal's ancestry, committed
or not. DRAIN, resends and ANCHOR votes use no quota.

\paragraph{2. Produce DATA while space remains.}
If $h(A_{\beta,r,g,0})<H_{\beta,r}$, generation $(\beta,r,g)$ is a
\defterm{DATA generation}. Each term begins at its anchor
$A_{\beta,r,g,t}$ and has
\[
 m_{\beta,r,g,t}=\min\{K,\;H_{\beta,r}-h(A_{\beta,r,g,t})\}
\]
new positions, numbered $j=0,\ldots,m_{\beta,r,g,t}-1$: at most $K$ per
term, and never more than the remaining quota. At position $j$, the DATA
block has
\[
 h=h(A_{\beta,r,g,t})+j+1,\qquad
 k=k_{\beta,r}(A_{\beta,r,g,t})+j+1.
\]
A new position is a \defterm{fresh position}, and its block is a
\defterm{fresh block}. The anchor itself is not fresh. The last position
$j=m_{\beta,r,g,t}-1$ is the \defterm{terminal position}, and its QC is
the \defterm{terminal QC}.

If the terminal DATA height of a term is below $H_{\beta,r}$, its QC
admits term $(\beta,r,g,t+1)$, with that terminal block as
$A_{\beta,r,g,t+1}$. If it equals $H_{\beta,r}$, there is no next term.
Only the last DATA term can have fewer than $K$ positions. The initial
generation is always a DATA generation, since $\TB_r\ge1$.

\paragraph{3. Use DRAIN when recovery starts at the endpoint.}
If $g>0$ and $h(A_{\beta,r,g,0})=H_{\beta,r}$, generation $(\beta,r,g)$ is
a \defterm{drain generation}. It has only term $t=0$ and position $j=0$,
with $m_{\beta,r,g,0}=1$. Its one fresh block is a DRAIN, whose parent is
$A_{\beta,r,g,0}$ and whose justification is the generation's ANCHOR QC.
It carries no requests, keeps DATA height $H_{\beta,r}$, and has no next
term.

DRAIN exists to supply a commit timer. ANCHOR voting starts no timer, so
re-certifying $A_{\beta,r,g,0}$ alone does not commit it. At height
$H_{\beta,r}$, a DATA child would exceed the quota. A DRAIN child is
voted on like any proposal, and its timer commits its ancestors,
including $T_{\beta,r}$.
\begin{itemize}
\item A DATA generation that reaches $H_{\beta,r}$ does not append DRAIN.
It waits for $T_{\beta,r}$'s timer or recovers after a quit.
\item A later drain generation may start from an earlier DRAIN.
Each interrupted recovery may add one DRAIN, so their number is not
bounded.
\item The choice depends on the recovery anchor's DATA height, not on
whether a replica has already committed $T_{\beta,r}$.
\end{itemize}
DRAIN stays in the regional log. The external chain contains only DATA,
and the next session inherits $T_{\beta,r}$.

\paragraph{Example.}
Let $h(A_{\beta,r,0,0})=100$, $\TB_r=10$, and $K=4$, so $H_{\beta,r}=110$.
\begin{center}\small
\begin{tabular}{@{}lll@{}}\toprule
View & Anchor height & Fresh DATA heights\\\midrule
$(\beta,r,0,0)$ & 100 & 101--104\\
$(\beta,r,0,1)$ & 104 & 105--108\\
$(\beta,r,0,2)$ & 108 & 109--110\\\bottomrule
\end{tabular}
\end{center}
Recovery from DATA 107 starts a new generation whose first term produces
108--110. Suppose instead that 110 has a QC but a Quit cancels its timer,
and recovery selects 110 again. DATA 111 would exceed the quota, and
ANCHOR voting on 110 starts no timer. So the recovery leader proposes a
DRAIN on top of 110; its timer commits 110. In both cases the endpoint
stays 110.

\subsection{Certificates and signed messages}
\label{sec:certificates}
\Cref{tab:certificates} lists the certificates of a session. Each needs
matching signatures from distinct members of $V_{i_{\beta,r}}$ with total
weight $>F_{i_{\beta,r}}$.

\begin{table}[htbp]\centering\small
\caption{Certificates of a session.}
\label{tab:certificates}
\begin{tabularx}{\linewidth}{@{}>{\raggedright\arraybackslash}p{.25\linewidth}>{\raggedright\arraybackslash}p{.24\linewidth}Y@{}}
\toprule
Certificate & Signed message & Purpose\\
\midrule
\defterm{Proposal QC} $\QC_{\beta,r,g,t}(B)$ & DATA-VOTE or DRAIN-VOTE & Justifies a
child. A terminal QC below $H_{\beta,r}$ also admits the next term.\\
\defterm{ANCHOR QC} & ANCHOR-VOTE & Re-certifies the opening $A_{\beta,r,g,0}$ in
recovery generation $g>0$, before its fresh proposals.\\
\defterm{BLAME proof} & BLAME & Freezes a generation (\cref{sec:quit}).\\
\defterm{Regional receipt} $\FC_{\beta,r}(B)$ & FINAL & Proves that a correct replica
committed DATA block $B$.\\
\bottomrule
\end{tabularx}
\end{table}

A correct replica casts at most one fresh vote per $(\beta,r,g,t,j)$ and
one ANCHOR vote per $(\beta,r,g)$. DATA and DRAIN votes follow the same
rule but are different signed types. No QC and no ANCHOR vote starts a
commit timer. A QC means only that enough votes exist; it does not mean
that any timer has expired.

\paragraph{Signed context.}
Every regional BFT signature binds the chain identifier, protocol domain, session
$(\beta,r)$, the hash of its fixed initial anchor $A_{\beta,r,0,0}$,
voting configuration, assigned region and message type. Binding the
anchor hash ties every signature to the exact start of its session. This
matters after a Force: lagging nodes may disagree about a discarded
suffix, and the hash stops evidence made under one start from being used
under another. View messages also bind $(g,t)$ and the authority-context
reference $\tagg{AUTH}_{\beta,r,g}$; their full view is $(\beta,r,g,t)$. The
other signed fields are:
\begin{itemize}
\item \textbf{PROPOSE and fresh votes:} the generation's opening
$A_{\beta,r,g,0}$, position $j$, block type, physical position, DATA
height and block hash;
\item \textbf{NEWVIEW and ANCHOR-VOTE:} the chosen opening. NEWVIEW
also binds its justification and entry proof;
\item \textbf{STATUS:} the sender's entry lock and entry proof;
\item \textbf{BLAME:} the view only, with no opening, block hash or
progress count. All complaints for a view sign the same statement;
\item \textbf{FINAL:} the session and DATA fields below, with no $g$ or $t$.
\end{itemize}
Displayed messages include $\beta$ and $r$ but omit the common chain,
protocol, region and configuration fields where they are clear from
context. A verifier also checks the boundary that assigned the pair
$(\beta,r)$; the pair alone is not evidence of authority.

\paragraph{Rank and the base justification.}
Within session $(\beta,r)$, a QC formed in term $(\beta,r,g,t)$ for block
$B$ has \defterm{rank} $(g,t,k_{\beta,r}(B))$, compared lexicographically.
The ANCHOR QC of recovery generation $g>0$ has rank $(g,0,k_{\beta,r}(A_{\beta,r,g,0}))$,
just below the generation's first fresh QC. A fixed total order on
certified content breaks ties; signer sets and encodings do not matter.
The \defterm{opening identity} of generation $g$ is
$(g,k_{\beta,r}(A_{\beta,r,g,0}),h(A_{\beta,r,g,0}),\Hash(A_{\beta,r,g,0}))$.
As in Part~I, ``opening'' always means the generation's opening block
$A_{\beta,r,g,0}$. QCs are ranked only within one session $(\beta,r)$.

The anchor $A_{\beta,r,0,0}$ also has a \defterm{base justification}
$Q_{\beta,r}^0$: the anchor together with its direct or covering
commitment evidence, carried by the handoff record that opened the
session, its \defterm{boundary witness} (\cref{sec:coverage,sec:witness}). Its justified block is exactly
$A_{\beta,r,0,0}$, even if the evidence uses a receipt for a higher
block of the old session; that higher block is evidence only.
$Q_{\beta,r}^0$ ranks below every QC of session $(\beta,r)$.
It directly justifies that session's first DATA and is not a new
certificate formed by its members. Recovery uses it when no QC of the
session is known. A QC of another session is
never a justification in session $(\beta,r)$, so it cannot replace
$A_{\beta,r,0,0}$.

\paragraph{Regional receipts.}
After a replica commits a DATA block $B$ of session $(\beta,r)$, by $B$'s
own timer or as an ancestor, it sends
\[
 \langle\tagg{FINAL},\beta,r,\mathit{cfg}_{\beta,r},h(B),\Hash(B)\rangle_v
\]
to its region. Matching FINAL messages of weight $>F_{i_{\beta,r}}$ form
the \defterm{regional receipt} $\FC_{\beta,r}(B)$. The receipt carries
$B$ and the ancestry needed to validate it, or refers to a validated
local cache.

On first obtaining a complete result, each correct holder offers it to
the Publish Layer and multicasts the full validated receipt, the DATA
and their dependencies to the fixed regional membership. Duplicates are
removed by session and DATA hash; the outgoing action is kept until it
is dispatched. This rule also covers a receipt used for Completion and
a holder that is Completed or archival. The holder may also send the
result directly to other regions for an early handoff.

\begin{definition}[\defterm{Valid regional receipt}]
\label{def:receipt}
A regional receipt $\FC_{\beta,r}(B)$ is valid when all of the following hold:
\begin{enumerate}[label=(\roman*)]
\item $B$ is a DATA block of session $(\beta,r)$, extends
$A_{\beta,r,0,0}$, and has $h(A_{\beta,r,0,0})<h(B)\le H_{\beta,r}$;
\item distinct members in session $(\beta,r)$'s fixed voting snapshot have
valid signatures on the same
$\langle\tagg{FINAL},\beta,r,\mathit{cfg}_{\beta,r},h(B),\Hash(B)\rangle$,
with total weight $>F_{i_{\beta,r}}$ and the signed context specified above; and
\item $B$ and its required ancestry pass the block-validation rules of
\cref{sec:admission}.
\end{enumerate}
Regional engines check the full evidence under \cref{sec:admission};
they do not replay other signers' timers. The global finality rule uses
\cref{sec:receipt-kernel} to establish the same intrinsic validity from
the exact DATA and FINAL signatures. The weight threshold guarantees a
correct signer, and that signer's own validation and commit ground the
conclusion.
\end{definition}
Here $\mathit{cfg}_{\beta,r}$ identifies session $(\beta,r)$'s \defterm{fixed
voting snapshot}. $\FC(B)$ without a session subscript means the receipt
from the session that produced $B$. An inherited anchor keeps that
receipt; it does not acquire a new one. The public checker of
\cref{sec:receipt-kernel} establishes condition~(iii) without requiring
any optional admission sidecar from a publisher. The regional engines
still check the full evidence.
\begin{itemize}
\item FINAL binds no generation or term, so FINAL messages for one block
combine across recoveries. The block hash keeps the block's signed
context.
\item If a DRAIN commits $T_{\beta,r}$, replicas send FINAL for
$T_{\beta,r}$, never for the DRAIN.
\item Replicas of session $(\beta,r)$ send no FINAL for $A_{\beta,r,0,0}$.
It keeps its original commitment evidence, which may come from an earlier
session if intermediate sessions added no DATA.
\item A Running replica may also commit $B$ and its uncommitted ancestors when it
obtains a valid $\FC_{\beta,r}(B)$ of its own session. Completion does
this for $T_{\beta,r}$ (\cref{alg:steady}). Such a commit sends no FINAL,
because the receipt already exists. Only Completion also cancels timers.
Receipts of other sessions serve only as history or anchor evidence;
they never change this session's locks or log.
\end{itemize}

\paragraph{Receipt coverage.}
\label{sec:coverage}
A valid $\FC_{\beta,r}(D)$ also proves commitment of every DATA ancestor
$B$ of $D$ in the same session: at least one correct signer committed
$D$, and committing $D$ commits its ancestors. The \defterm{covering
evidence} for $B$ is the valid receipt for $D$ plus the fully verified
ancestry from $D$ to $B$. It does not claim that the signatures name $B$
directly. For example, a receipt for 110 with the exact path back to 109
proves that 109 was committed.

Coverage never crosses sessions: a receipt of B's session for 111 does
not replace A's evidence for 110. Coverage also does not show that $B$
belongs to the authorized external branch. A later Force may discard $D$
while keeping $B$; the regional evidence for $B$ still holds. External
observers, retained-anchor witnesses and base justifications may use
direct or covering evidence, with no extra FINAL messages or certificate
round. We write ``\defterm{commitment evidence} for $B$'' for either form.

\subsection{Admission, validation, and engine states}
\label{sec:admission}
\paragraph{\defterm{Admission.}}
A replica acts in a view only after it \emph{admits} that view.
Replica $v$ admits views as follows.
\begin{itemize}
\item $(\beta,r,0,0)$ is admitted at $v$'s initial Enter for the boundary
that opens session $(\beta,r)$ (\cref{def:ready}). The boundary witness is
the \defterm{entry proof}. It fixes $A_{\beta,r,0,0}$ and
$Q^0_{\beta,r}$ directly; generation zero has no NEWVIEW or ANCHOR vote.
\item $(\beta,r,g,0)$ with $g>0$ is admitted once $v$ holds
$\tagg{AUTH}_{\beta,r,g}$ and a quit proof for generation $g-1$. That proof is
the \defterm{entry proof}. The NEWVIEW is justified by a QC of an earlier
generation of session $(\beta,r)$, or by $Q_{\beta,r}^0$.
\item $(\beta,r,g,t)$ with $t>0$ is admitted once $v$ also holds the
terminal QCs of terms $0,\ldots,t-1$, each for DATA below $H_{\beta,r}$.
\end{itemize}
Term zero is admitted at entry, so its progress checks run even if
the leader is silent. The first proposal of generation zero uses
$Q^0_{\beta,r}$; recovery generations require their ANCHOR QC.
Subsequent proposals follow \cref{tab:parent}. A view message is
\defterm{admissible} if its view is admitted, its authority context matches
$(\beta,r,g)$ and the fixed session configuration, and its signatures are
valid. Leader-signed messages are checked against $L(\beta,r,g,t)$ for
that exact full view.
FINAL instead uses the initial session context of \cref{sec:receipt-kernel},
with no generation or term binding.

A verifier that checks a remote session's full regional history has no
local admission event for it. There, a view counts as admitted if the
evidence contains its \defterm{admission witness}: the boundary witness
for $(\beta,r,0,0)$, the chain of entry proofs for $(\beta,r,g,0)$, and
the earlier terminal QCs for $t>0$.

\paragraph{\defterm{Complete regional validation.}}
Complete regional validation checks signatures, authority, admission,
block type, quota, both positions, the L1 and Espresso reference rules,
the parent rule (\cref{tab:parent}), parent justifications, hash links and
application validity throughout the required ancestry. Every relevant
ancestor comes with its signed proposal and block data; a chain of
certificate hashes alone is not enough. Dependencies come from the
message or from a validated local cache. Correct senders send
\defterm{complete validation bundles}. A reference to a dependency is
allowed only if the receiver can resolve it within the delay bound of
\cref{ass:network}. Bundles list each dependency once.

Regional Vote, recovery, and processing of own-session evidence and
receipts always use this complete procedure. Ready also checks the
target's admission, exact state and boundary inputs in full; only the
receipt of the inherited anchor, which another session produced, may be
checked with the public kernel of \cref{sec:receipt-kernel}.

\paragraph{Evidence takes priority.}
Before validating a message or discarding a duplicate, inspect all signed
material in it for a quit proof (\cref{sec:quit}). A proof needs its own
valid signatures and admission witnesses. A valid proof counts even if
the message that carries it is invalid or incomplete. Compare incoming
material with all stored material, not only with objects inside the
message.

\paragraph{Processing proposals.}
While the program is Running, every admitted fresh proposal of its
current generation that $v$ has not voted for goes through the Vote step.
This includes proposals inside attached histories, proposals
of earlier terms, and proposals whose blocks are already committed as
ancestors. A proposal's identity is $(\beta,r,g,t,j,\Hash(B))$; the signer
set or encoding of its parent QC does not change it. Duplicates are
removed by this identity after evidence inspection. Old-generation
proposals are stored without new votes or timers. Future-generation
messages wait for entry. After termination the Manager stores late
evidence without restarting BFT.

\paragraph{Engine states.}
A replica runs one engine for each session pair $(\beta,r)$ assigned to
its region. Engines of different sessions stay separate, even if they
share a turn index. The engine of session $(\beta,r)$ has three core
phases. Its host program is Running while the engine is Active or
Frozen. Completion or Force-close terminates the program
(\cref{sec:session-lifecycle}); the Manager keeps its required data.
\begin{itemize}
\item \textbf{\defterm{Active}:} the current generation is admitted and
unfrozen. The engine proposes, votes and changes terms.
\item \textbf{\defterm{Frozen}:} the current generation has a quit proof.
The engine stops proposals, votes and term changes. It still processes
evidence, old votes, QCs, receipts and recovery.
\item \textbf{\defterm{Completed}:} the replica holds
$\FC_{\beta,r}(T_{\beta,r})$ satisfying \cref{def:receipt}. Production in
session $(\beta,r)$ is over, even if the external handoff is still
pending. The Manager continues the required data service.
\end{itemize}
Committing $T_{\beta,r}$ locally does not complete the engine: other
members may still need this replica's recovery votes before a receipt
forms. Each message goes to the engine named by its signed pair
$(\beta,r)$, so evidence from one session or branch never freezes
another.

These states describe the local engine, not external finality.
A node that has not learned an effective Force may still run the old
engine. Once its Manager verifies the Force, it terminates every
old-version program and cancels its timers (\cref{sec:force-close}).
Existing DATA and evidence remain in storage.

\begin{table}[htbp]\centering\small
\caption{Parent selection for the next fresh position $j$ of term $(\beta,r,g,t)$.}
\label{tab:parent}
\begin{tabularx}{\linewidth}{@{}l Y Y@{}}
\toprule
Condition & $B_{\mathrm{parent}}$ & $Q_{\mathrm{parent}}$\\
\midrule
$g=t=j=0$ & Fixed anchor $A_{\beta,r,0,0}$ & Base justification $Q^0_{\beta,r}$\\
$g>0,\ t=j=0$ & Recovery opening $A_{\beta,r,g,0}$ & Its ANCHOR QC in $(\beta,r,g,0)$\\
$j=0,t>0$ & Anchor $A_{\beta,r,g,t}$, the terminal block of $(\beta,r,g,t-1)$ & Its proposal QC in $(\beta,r,g,t-1)$\\
$1\le j<m_{\beta,r,g,t}$ & Position $j-1$ of $(\beta,r,g,t)$ & Its proposal QC in $(\beta,r,g,t)$\\
\bottomrule
\end{tabularx}
\end{table}

\section{Regional protocols}
\label{sec:protocols}
\Cref{alg:steady,alg:recovery} run inside the engine of one session
$(\beta,r)$. They are Protocols 1 and 2 of Part~I~\cite{part1}, with the
changes listed below. \RR{Red text} marks where they differ from Part~I.
The Espresso reference $\xi$ is defined in \cref{sec:blocks}.

The life of a session is short to state. It starts from its fixed
handoff under \cref{alg:bootstrap} (\cref{sec:ready}). It then runs
\cref{alg:steady}, and \cref{alg:recovery} whenever a generation is
frozen. It stops normally through Completion, when it holds a receipt
for its terminal DATA, or immediately when its Manager verifies a Force
closing its version. Quit cancels generation timers; Completion and
Force-close cancel all remaining session timers.

\paragraph{Notation.}
The full view is $(\beta,r,g,t)$ throughout. Every anchor, receipt and
entry time below names its branch and session explicitly. In these two
protocols, $\Delta$ and $F$ mean $\Delta_{i_{\beta,r}}$ and
$F_{i_{\beta,r}}$. The voting configuration $\mathit{cfg}_{\beta,r}$ stays
fixed within the session; $\tagg{AUTH}_{\beta,r,g}$ fixes its generation's
leader rule, evaluated at term $t$. A later branch reusing $r$ has a new
authority context and may select a different snapshot. Only $\TB_r$
remains shared by that turn index.

\paragraph{Changes from Part~I.}
\begin{enumerate}[label=\arabic*.]
\item \emph{Stake.} Certificates need weight $>F$, not $f+1$ signatures.
\item \emph{Blocks.} A block also carries its session, physical
position, type, L1 reference $\rho$, and Espresso reference $\xi$
(\cref{sec:blocks}). Validity checks all of them.
\item \emph{Quota.} A term has $m_{\beta,r,g,t}\le K$ positions. No term
follows $H_{\beta,r}$.
\item \emph{DRAIN.} A drain generation proposes one DRAIN. Its vote
counts as progress.
\item \emph{Receipts.} Replicas send FINAL after committing DATA. A
receipt for $T_{\beta,r}$ completes the engine (step~5 of \cref{alg:steady}).
\item \emph{Timers.} Completion and verified Force-close cancel all
remaining session timers.
\item \emph{Recovery entry.} After the $\Delta$ wait, a replica enters
when the finalized recovery record makes $\tagg{AUTH}_{\beta,r,g+1}$
available. $Q_{\beta,r}^0$ can serve as the
entry lock and as the justification.
\item \emph{Initial entry.} Every new session, including genesis, starts
directly from its fixed anchor and $Q^0_{\beta,r}$, without an ANCHOR QC.
\item \emph{Sessions.} Each message goes to the engine of its session.
The initial session $(\beta,0)$ also opens through \cref{alg:bootstrap}.
\end{enumerate}

\paragraph{Program guard.}
Every BFT action requires a Running program whose version is not marked
closed: proposals, votes, STATUS, NEWVIEW, BLAME, commit-timer callbacks,
term changes and recovery. The same guard applies to buffered messages
and authority-arrival callbacks. A verified Force closes the guard and
cancels all outstanding BFT work before the next action. The Manager may
still store and serve existing data and evidence; it issues no new BFT
signature for the terminated program. Program events never modify $G$.

\paragraph{Message-processing order.}\leavevmode
\RR{Replica $v$ first passes each message to the engine of the session
pair $(\beta,r)$ that the message binds. L1 updates go to the Force handler in Part~IV~\cite{part4}. Boundary witnesses and receipts of
other sessions go to \cref{alg:rotation,alg:global-rule}.} The engine handles each received
message in this order:
\begin{enumerate}[label=(\arabic*)]
\item check all signed material for a quit proof;
\item check completeness and validity;
\item process attached DATA \RR{and DRAIN} proposals under Vote, parents
first;
\item process attached terminal QCs under step~4 of \cref{alg:steady},
in increasing term order;
\item process the message itself.
\end{enumerate}
When a message and a timer fall at the same time, $v$ handles the
message first.

\subsection{Steady state}
\label{sec:steady}
\Cref{alg:steady} gives the steady state of a session: proposals, votes,
timers, term rotation, and completion.

\begin{protocol}{alg:steady}{Propose, vote, commit, normal rotation\RR{, and completion}}
Replica $v$ runs these steps \RR{in the engine of session $(\beta,r)$}. Its
current term is $(\beta,r,g,t)$.
\begin{enumerate}[label=\textbf{\arabic*)},itemsep=3pt]
\item \textbf{Propose.} Run this step only if $v=L(\beta,r,g,t)$, $g$ is
unfrozen\RR{, and session $(\beta,r)$ is not complete}.
\RR{In a DATA generation,} for each position
$j=0,\ldots,\RR{m_{\beta,r,g,t}}-1$, in order:
 \begin{itemize}
 \item Wait until $v$ holds $B_{\mathrm{parent}}$ and
       $Q_{\mathrm{parent}}$ for $j$ (\cref{tab:parent}).
 \item Build \RR{$B=(\beta,r,k,h,\tagg{DATA},x,\Hash(B_{\mathrm{parent}}),\rho,\xi)$,
       with $k$ and $h$ from \cref{sec:quota} and the references
$\rho,\xi$ from \cref{sec:blocks}}.
       The batch $x$ is an opaque application payload.
 \item Broadcast
       $P=\langle\tagg{PROPOSE},\beta,r,g,t,j,B,Q_{\mathrm{parent}}\rangle_v$
       with its validation bundle \RR{to regional replicas. }
 \item Process $P$ under Vote.
 \end{itemize}
 \RR{In a drain generation, propose once, at $(t,j)=(0,0)$. The block is
 a DRAIN with parent $A_{\beta,r,g,0}$, no batch, DATA height $H_{\beta,r}$, and the L1
 and Espresso references of $T_{\beta,r}$. Broadcast and process it as above.}
 Do not wait for the parent to commit. Never sign two different
 blocks for one position. A resend repeats the same $P$.

\item \textbf{Vote.} Run this step on every DATA \RR{or DRAIN} proposal
$P$ of term $(\beta,r,g_P,t_P)$ at position $j_P$. Here $t_P$ may be earlier
than $t$. The step also applies to proposals inside attached histories.
 \begin{enumerate}[label=(\alph*)]
 \item If $P$ or its attached material contains a quit proof for $g$,
       run Quit and stop.
 \item If $g_P<g$, keep $P$ and stop. If $g_P>g$, buffer $P$
       until $v$ enters generation $g_P$.
 \item If $g$ is frozen \RR{or session $(\beta,r)$ is complete}, keep $P$ and
       stop.
 \item If a dependency is missing, keep $P$, request the dependency,
       and stop. Rerun this step when it arrives.
 \item If $P$ is invalid, discard it and stop. \RR{Validity includes the
       block type, the quota, both positions, and the L1 and Espresso references.}
 \item If $v$ has already voted for $P$, stop.
       A duplicate never restarts a timer.
 \item Otherwise, let $\tau_v(P)$ be the current time. Forward $P$ and
       its bundle to all \RR{regional replicas. }
       Broadcast
       $\langle\tagg{DATA\mbox{-}VOTE},\beta,r,g,t_P,j_P,\Hash(B)\rangle_v$
       \RR{(DRAIN-VOTE if $B$ is a DRAIN)}.
       Start the commit timer of $P$. It expires at $\tau_v(P)+2\Delta$.
 \end{enumerate}

\item \textbf{Commit.} When the timer of $P$ expires, commit $B$ and
its uncommitted ancestors. \RR{For each DATA block $D$ of session $(\beta,r)$
that this commits, send
$\langle\tagg{FINAL},\beta,r,\mathit{cfg}_{\beta,r},h(D),\Hash(D)\rangle_v$ to the regional fixed membership and the Publish Layer collection interface. A committed DRAIN yields FINAL for $T_{\beta,r}$, never for
itself. } Quit \RR{, Completion and verified Force-close} cancel timers.
Entering a new term does not cancel one. A QC never starts a timer.

\item \textbf{Terminal QC and rotation.} Run this step the first time
$v$ obtains a valid terminal QC for a term $(\beta,r,g,t')$.
``Obtains'' means that $v$ assembles the QC from votes or finds it in
any message, including as another proposal's $Q_{\mathrm{parent}}$.
 \begin{enumerate}[label=(\alph*)]
 \item Process each unseen proposal in the QC's history under Vote,
       parents first.
 \item If $g$ is now frozen, stop. \RR{If the QC certifies DATA at
       $H_{\beta,r}$, or a DRAIN, stop: no next term exists.}
 \item Send the QC and its validation bundle to $L(\beta,r,g,t'+1)$.
 \item If $t\le t'$, enter $(\beta,r,g,t'+1)$ and set $e_v(\beta,r,g,t'+1)$ to the
       current time. Cancel the progress checks of all earlier
       terms.
 \end{enumerate}
 A QC never moves $v$ to an earlier term. It never changes a recorded
 entry time.

\item \RR{\textbf{Completion.} Run this step the first time $v$ obtains a
regional receipt $\FC_{\beta,r}(T_{\beta,r})$ satisfying \cref{def:receipt}.}
 \begin{itemize}\color{rrRed}
 \item Multicast the complete validated terminal receipt and its dependencies
       to the fixed regional membership, once, and persist this outgoing action.
 \item Commit $T_{\beta,r}$ and its uncommitted ancestors.
 \item Mark session $(\beta,r)$ Completed. Stop its proposals, votes, progress
       checks, term changes, and recovery. Cancel its unexpired timers.
 \item Offer the complete terminal result to the Publish Layer.
       It may also be sent directly to the next region.
 \item Keep serving data. \Cref{alg:rotation} runs the handoff.
 \end{itemize}
 \RR{A force instruction never runs this step.}
\end{enumerate}
\end{protocol}

\Needspace{9\baselineskip}
The \defterm{acceptance time} $\tau_v(P)$ is when complete validation and the first
vote occur, not an earlier arrival of an incomplete packet. Each replica
records its own entry time for each term. In session $(\beta,r)$, $e_v(\beta,r,0,0)$ is
set at the atomic Enter call, in step~1 of \cref{alg:bootstrap}. For $g>0$, $e_v(\beta,r,g,0)$ is
set in step~3 of \cref{alg:recovery}. For $t>0$, $e_v(\beta,r,g,t)$ is set in
step~4 of \cref{alg:steady}. A term's progress checks therefore start
when the term is admitted locally, not when its leader's first proposal
arrives. A silent new leader cannot delay them.

A proposal or QC at height $H_{\beta,r}$ does not end the session. Without Quit or Force-close, the timer of $T_{\beta,r}$ commits it;
no later block is needed. Only a receipt for $T_{\beta,r}$ completes the
engine normally.

\subsection{Quit proofs and recovery}
\label{sec:quit}
A quit proof freezes a generation. \Cref{alg:recovery} then moves the
session to the next generation.

\begin{protocol}{alg:recovery}{Blame, quit, and recovery}
Replica $v$ runs these steps \RR{in the engine of session $(\beta,r)$}. Its
current term is $(\beta,r,g,t)$.
\begin{enumerate}[label=\textbf{\arabic*)},itemsep=3pt]
\item \textbf{Blame.} For each $p=1,\ldots,\RR{m_{\beta,r,g,t}}$ (one
\RR{in a drain generation}), run a
progress check at time
\[
 e_v(\beta,r,g,t)+(2p+4)\Delta.
\]
\RR{Schedule the check for $p=1$ at entry, even if the opening is not yet
known. Schedule the others once $m_{\beta,r,g,t}$ is known. A check whose time has
already passed runs at once.}
At that time, broadcast $\langle\tagg{BLAME},\beta,r,g,t\rangle_v$ if all of
these hold:
 \begin{itemize}
 \item $v$'s current term is still $(\beta,r,g,t)$, $g$ is unfrozen\RR{, and
       session $(\beta,r)$ is not complete};
 \item $v$ has cast DATA \RR{or DRAIN} votes for fewer than $p$ positions
       of term $(\beta,r,g,t)$ (ANCHOR votes do not count);
 \item $v$ has not yet sent BLAME for $(\beta,r,g,t)$.
 \end{itemize}

\item \textbf{Quit.} Run this step the first time $v$ holds a quit proof
for its current generation $g$.
 \begin{itemize}
 \item Freeze $g$. Stop all progress checks, proposals, votes, and
       term changes in $g$. Cancel every unexpired commit timer of $g$.
 \item Set $a_v$ to the current time. Broadcast the proof.
 \item Keep all commitments. Keep collecting votes and QCs of $g$.
       They can raise the entry lock chosen in step~3. \RR{Keep
       processing regional receipts (step~5 of \cref{alg:steady}).}
 \item If the proof is of type (b) or (c), submit it to L1. Do
       not wait for L1 to process it.
 \end{itemize}

\item \textbf{Status and entry.} \RR{Skip this step if session $(\beta,r)$ is
complete.} \RR{Wait until at least $a_v+\Delta$ and until
the recovery record for $(\beta,r,g+1)$ is finalized and its authority
context is available. Set $e_v(\beta,r,g+1,0)$ to the time both conditions
hold, and derive $\tagg{AUTH}_{\beta,r,g+1}$ from that record and the
session's initial authority.} Then:
 \begin{itemize}
 \item Set the entry lock $Q_v$ to the highest-ranked QC
       (DATA\RR{, DRAIN,} or ANCHOR) that $v$ knows from generations
       $0,\ldots,g$\RR{, or to $Q_{\beta,r}^0$ if it knows none}.
       $Q_v$ never changes after this. QCs that arrive later do not
       update it.
 \item Enter $(\beta,r,g+1,0)$ as unfrozen and start its progress checks.
       \RR{Keep $A_{\beta,r,0,0}$, $H_{\beta,r}$, and the voting snapshot of session $(\beta,r)$.}
 \item Send $\langle\tagg{STATUS},\beta,r,g+1,0,Q_v,E\rangle_v$ and its
       validation bundle to $L(\beta,r,g+1,0)$. Here $E$ is the quit proof for
       $g$, which is also the entry proof of $g+1$.
 \item Process any buffered proofs for $g+1$.
 \end{itemize}

\item \textbf{New-view.} Run this step only at $L(\beta,r,g+1,0)$. Wait $2\Delta$
after its own entry. If $g+1$ is still current and unfrozen:
 \begin{itemize}
 \item Let $Q^*$ be the highest-ranked QC from generations $0,\ldots,g$
       that the leader knows, including the QCs in STATUS
       messages\RR{; if it knows none, let $Q^*=Q_{\beta,r}^0$}.
 \item Let $A_{\beta,r,g+1,0}$ be the block that $Q^*$ certifies. \RR{If
       $h(A_{\beta,r,g+1,0})<H_{\beta,r}$, generation $g+1$ is a DATA generation.
       Otherwise it is a drain generation.}
 \item Broadcast
       $\langle\tagg{NEWVIEW},\beta,r,g+1,0,A_{\beta,r,g+1,0},Q^*,E\rangle_{L(\beta,r,g+1,0)}$
       with its validation bundle. Choose only one opening block for
       $g+1$.
 \end{itemize}

\item \textbf{Anchor vote.} On receiving a NEWVIEW for $g+1$:
 \begin{enumerate}[label=(\alph*)]
 \item If it contains a quit proof for $g+1$, run Quit and stop.
 \item Validate its leader, entry proof $E$, opening block $A_{\beta,r,g+1,0}$,
       $Q^*$, and dependencies. If data are missing, request them and
       wait without voting. If the message is invalid, discard it.
 \item Vote only if all of these hold: $g+1$ is current and unfrozen; $v$
       has cast no ANCHOR vote in $g+1$; and $\rank(Q^*)\ge\rank(Q_v)$.
 \item To vote, forward the NEWVIEW to all replicas. Broadcast
       \[\langle\tagg{ANCHOR\mbox{-}VOTE},\beta,r,g+1,0,\Hash(A_{\beta,r,g+1,0})\rangle_v.\]
       Start no timer. Use no fresh position.
 \end{enumerate}
 Matching ANCHOR votes \RR{of total weight $>F$} form the ANCHOR QC.
 When $L(\beta,r,g+1,0)$ obtains it, the leader starts Propose in
 \cref{alg:steady} at $j=0$\RR{; in a drain generation, this is the
 DRAIN}. Every replica, including one that cast no ANCHOR vote,
 processes DATA \RR{and DRAIN} proposals backed by a valid ANCHOR QC
 under Vote.
\end{enumerate}
\end{protocol}

\paragraph{\defterm{Complete quit proofs.}}
For an admitted generation $g$ \RR{of session $(\beta,r)$}, a proof consists of
either:
\begin{enumerate}[label=(\alph*)]
\item distinct signatures \RR{of total weight $>F$} on the same
$\langle\tagg{BLAME},\beta,r,g,t\rangle$ for an admitted term, with that term's
admission witness;
\item two different admissible DATA \RR{or DRAIN} blocks signed by
$L(\beta,r,g,t)$ at the same $(\beta,r,g,t,j)$, with their admission \RR{and quota}
witnesses; or
\item two valid NEWVIEWs signed by $L(\beta,r,g,0)$ selecting different opening
blocks for a recovery generation $g>0$.
\end{enumerate}
\RR{A DATA block beyond the quota, or a DRAIN outside a drain generation,
is invalid. It is rejected and never counts for (b).}
Different openings count even when their blocks are comparable, because
they fix different starting heights $h(A_{\beta,r,g,0})$. The height formula then
maps the same position to different heights.

Different QCs for the same opening block are not equivocation. Evidence
from any admitted term can freeze its generation, including a term not
yet locally reached. Old-generation proofs are archival and may be
submitted to L1. \RR{Evidence from another session never freezes this
engine.} Future-generation proofs are buffered until entry. Neither
resets an unrelated recovery.

The Publish Layer uses these DATA and FINAL receipts. No publication
acknowledgement is required for regional commitment or Completion.
Participants may forward regional bundles directly; no external send
or Espresso confirmation delays a regional vote or timer.

The entry-lock test governs ANCHOR voting only. A replica that has
committed $T_{\beta,r}$ locally still takes part in recovery until it holds
$\FC_{\beta,r}(T_{\beta,r})$, because other members may need its votes. This continued-participation
rule applies until Completion or a verified Force terminates the program.

\section{Regional and global finality}
\label{sec:external}
Regional commitment and global acceptance are separate events. A node
follows a production branch $R$ and a continuous global DATA prefix
$G$. The branch may extend several sessions beyond $G$. A Force
closes that extension, never $G$ itself.

This section first states what a regional receipt proves
(\cref{sec:regional-finality}) and gives the public check of a receipt
(\cref{sec:receipt-kernel}). It then defines the global chain $G$ and
global finality (\cref{sec:global-chain}), and the rule by which every
verifier computes $G$ from the Espresso log and L1
(\cref{sec:exports}).

\subsection{Regional finality}
\label{sec:regional-finality}
\begin{definition}[\defterm{Regional finality}]
\label{def:regional-finality}
A DATA block $B$ has regional finality on a specified production branch
when its boundary history assigns it to its claimed region and session,
it has valid direct or same-session covering commitment evidence, and
the DATA path to it is valid and continuous. Each session on that path
supplies its own authority and commitment evidence.
\end{definition}
The region is derived from the boundary before its receipt is checked.
A valid receipt therefore cannot authorize another region's DATA or a
different session. Current live/off state does not replace historical
authority. When a Force closes an unconfirmed suffix, its receipts remain
evidence about the old regional engine, but they no longer let that
suffix extend the global prefix.

The regional agreement and progress statements below concern a branch
before its Force activation, while the stated regional-service premises
hold. Force-close can cancel unfinished work. The discarded suffix has
no post-Force regional agreement or completion guarantee
(\cref{sec:force-close}); global acceptance still preserves $G$.


\subsection{Public receipt verification}
\label{sec:receipt-kernel}
This subsection answers one question: given a DATA block $B$, can anyone
check from public data that $B$ is regionally final
(\cref{def:regional-finality})? The answer is the check
$\tagg{CheckFINAL}$. It needs three things: the context of the session
that produced $B$, which the verifier derives itself; FINAL signatures
on a root block $D$; and the exact bytes that link $D$ back to $B$. It
does not need the proposal or recovery history that a publisher may
attach, which can name recovery authorities the verifier does not have.


\paragraph{The session context $C$.}
Every check is made against the \defterm{session context} $C$ of the
producing session $(\beta,r)$. $C$ consists of two authenticated
objects: the session's incoming boundary (the handoff record that opened
it, \cref{sec:witness}) and its initial authority $\tagg{AUTH}_{\beta,r,0}$.
Both are checked against the producing session, including its historical
configuration; a newer branch's authority cannot replace them.
They fix the following fields.
\begin{center}\small
\begin{tabularx}{\linewidth}{@{}>{\raggedright\arraybackslash}p{.36\linewidth}>{\raggedright\arraybackslash}p{.2\linewidth}Y@{}}\toprule
Field & Source & Example: A's $(7,20)$\\\midrule
chain identifier and protocol domain & boundary & this chain\\
producing pair $(\beta,r)$ & boundary & $(7,20)$\\
assigned region $i$ & boundary & A\\
anchor $A_{\beta,r,0,0}$: height, hash, original producer & boundary & DATA 100\\
fixed voting snapshot $\mathit{cfg}_{\beta,r}$: members, weights, $F_i$ & $\tagg{AUTH}_{\beta,r,0}$ & $\mathit{cfg}_{7,20}$\\
DATA quota $\TB_r$, hence $H_{\beta,r}$ & $\tagg{AUTH}_{\beta,r,0}$ & 10, so $H=110$\\\bottomrule
\end{tabularx}
\end{center}
The verifier derives $C$ from public history, never from the sender of a
record. The global finality rule (\cref{alg:global-rule}) installs a
regular boundary when it accepts the predecessor's terminal DATA, a
forced boundary at a Force, and fixes the region by applying
$\tagg{NextLive}$ to the branch's immutable map; the
initial checkpoint supplies the genesis boundary; a handoff verifier
derives the same boundary with $\tagg{VerifyHandoffHistory}$
(\cref{sec:verify-handoff}). $C$ contains no recovery authority
$\tagg{AUTH}_{\beta,r,g}$ with $g>0$ and no generation or term: FINAL does not
bind them, and the voting snapshot is the same in every generation of
the session.

\paragraph{The check.}
Write
\[
 \tagg{CheckFINAL}(C,B,D,\Sigma,\pi).
\]
Its inputs are the following.
\begin{center}\small
\begin{tabularx}{\linewidth}{@{}>{\raggedright\arraybackslash}p{.05\linewidth}Y>{\raggedright\arraybackslash}p{.15\linewidth}>{\raggedright\arraybackslash}p{.22\linewidth}@{}}\toprule
Input & What it is & Supplied by & Example: A's $(7,20)$\\\midrule
$C$ & the session context above & derived by the verifier & anchor DATA 100, quota 10\\
$B$ & the candidate DATA to check & publisher & DATA 103\\
$D$ & the root DATA that carries the FINAL certificate: $B$ itself, or a same-session descendant of $B$ & publisher & DATA 105\\
$\Sigma$ & the FINAL signatures on $D$ & publisher & FINAL on DATA 105\\
$\pi$ & the exact hash-linked path from $D$ back to $B$, with full bodies & publisher & $105\to104\to103$\\\bottomrule
\end{tabularx}
\end{center}
The triple $(D,\Sigma,\pi)$, read against $C$, forms the \defterm{kernel} for $B$.
The check has two parts.
\begin{enumerate}[label=(\arabic*),itemsep=2pt]
\item \textbf{Signatures.} Write $\tagg{RawFINAL}_C(D,\Sigma)$ when the
bytes $D$ have the named hash and height, and $\Sigma$ contains valid
FINAL signatures from distinct members of $\mathit{cfg}_{\beta,r}$, with total
weight $>F_i$. The signed context must match $C$ in full: chain,
protocol domain, message type, producing pair, initial anchor hash,
assigned region and configuration, even where the displayed FINAL
notation omits those fields. ``Raw'' means a cryptographic and type check only. It does not
validate the block, include \cref{def:receipt}(iii), read a generation
or term, or check any attached admission witness.
\item \textbf{Path.} Each body on $\pi$ is an actual DATA or DRAIN body
with the correct hash, parent link and position; $\pi$ ends at $B$; $D$
and $B$ belong to the same session $(\beta,r)$; and
$h(B)\le H_{\beta,r}$. For a direct receipt, $D=B$ and $\pi$ is just $B$.
A hash without its payload is not enough, and a later session's receipt
never covers an earlier session's DATA.
\end{enumerate}
The check returns \emph{true} if both parts pass, \emph{false} if one
fails, and \emph{incomplete} if a required object is absent. The caller
handles \emph{incomplete} under the missing-object rule; it is not an
instruction to fetch optional sidecars. Required bytes must be supplied
or resolved from established public objects.

In the example, part (1) checks that the FINAL signatures on DATA 105
have weight $>F_A$ under $\mathit{cfg}_{7,20}$. Part (2) checks that
$105\to104\to103$ is correctly linked inside $(7,20)$ and that
$103\le110$. The check returns \emph{true}: DATA 103 is regionally final.



\subsection{The global chain}
\label{sec:global-finality}
\label{sec:global-chain}
The \defterm{global chain} $G$ is the sequence of DATA blocks that every
verifier accepts. It starts at the genesis anchor $A_0$ and passes
through a sequence of sessions.

\paragraph{Session descriptors and block-bound records.}
To accept a DATA block, a verifier must know which session produced it
and how that session was opened. A session descriptor holds exactly this
information, and every accepted block keeps the descriptor it was
accepted under. A \defterm{session descriptor} is
\[
 \nu=(\beta,\ r,\ i,\ A,\ H,\ q^{\rm open},\ \mathit{cause}),\qquad
 \mathit{cause}\in\{\tagg{GENESIS},\ \tagg{REGULAR},\
 \tagg{FORCE}(c_f,q_f)\}.
\]
Its session identifier is $\operatorname{sid}(\nu)=(\beta,r)$; $i$ is
the selected region; $A$ is the inherited anchor;
$H=h(A)+\TB_r$ is the terminal DATA height; and $q^{\rm open}$ is the
lower bound on the Espresso header that the session's first DATA may
reference. It is not a selection time. The cause says how the
session was opened; for a Force, $c_f$ identifies the executing atomic
update and $q_f$ is its activation block.  We write $\beta(\nu)$,
$r(\nu)$, $i(\nu)$,
$A(\nu)$, $H(\nu)$, $q^{\rm open}(\nu)$ and
$\mathit{cause}(\nu)$ for these fields.  The descriptor fixes the session context $C_\nu$
of \cref{sec:receipt-kernel}: $\nu$ supplies the pair, region and anchor,
and $\mathit{cfg}_{\beta,r}$ and $\TB_r$ come from $\tagg{AUTH}_{\beta,r,0}$.  The L1
bound $\lambda^{\rm open}_{\beta,r}$ is not part of $\nu$; only regional
validation uses it.

Whenever the global finality rule appends a non-genesis DATA block
$B$ to $G$, it binds the immutable \defterm{session record}
\[
 \sigma_G(B):=\nu
\]
using the complete descriptor under which $B$ was accepted. This is
global-chain metadata, not a separately submitted or signed L1 or
Espresso record.  Every DATA produced by the same fixed session binds
the same descriptor value.  A selected session that contributes no DATA
to $G$ leaves no session record in $G$.

\begin{definition}[\defterm{Global chain}]
\label{def:global-chain}
A global chain is a sequence
\[
 G=A_0\cdot B_1\cdots B_n
\]
together with the map $B\mapsto\sigma_G(B)$ on its non-genesis blocks,
such that:
\begin{enumerate}[label=(\roman*),itemsep=2pt]
\item \textbf{Continuity.} Every $B_j$ has $B_{j-1}$ as its parent and
the next DATA height, where $B_0=A_0$.
\item \textbf{Ownership.} If $\sigma_G(B_j)=\nu_j$, then $B_j$ is
produced by $\operatorname{sid}(\nu_j)$, has height at most $H(\nu_j)$,
and is regionally final under $C_{\nu_j}$
(\cref{def:regional-finality}).
\item \textbf{Session runs.} Partition $B_1,\ldots,B_n$ into maximal
nonempty consecutive runs on which $\sigma_G$ is constant.  The first
run, if any, has turn index $0$ and anchor $A_0$.  If a run with
descriptor $\nu'$ follows one with descriptor $\nu$, then
$A(\nu')$ is the last block of the earlier run and
$r(\nu')=r(\nu)+1$.  If
$\mathit{cause}(\nu')=\tagg{REGULAR}$, the earlier run ends at
$H(\nu)$ and $\beta(\nu')=\beta(\nu)$.  If
$\mathit{cause}(\nu')=\tagg{FORCE}(c_f,q_f)$, then
$\beta(\nu')>\beta(\nu)$.  A run that ends below its terminal height can
only be followed by a Force-opened run.
\end{enumerate}
\end{definition}
No empty run appears in $G$.  A selected session with no accepted DATA
is represented by $\tagg{NEXT}(\nu)$ only while it owns the next block;
if a Force replaces it, it leaves no block or session record in $G$.
Repeated Forces before the next accepted DATA may therefore make the
branch version jump by more than one between two adjacent runs.

The rule of \cref{sec:exports} fixes which global chain each verifier
holds at each Espresso height. In the running example, each DATA in the
listed range binds the displayed record:
\begin{center}\small
\begin{tabular}{@{}ll@{}}\toprule
DATA in $G$ & Bound session record\\\midrule
$101,\ldots,108$ & $(7,20,\text{A},100,110,q_{\rm A},\tagg{REGULAR})$\\
$109',110',\ldots$ & $(8,21,\text{B},108,108+\TB_{21},q_f,\tagg{FORCE}(c_f,q_f))$\\\bottomrule
\end{tabular}
\end{center}
A's run stops at 108 because the Force closed version 7 before 109
was accepted. B's regular session $(7,21)$ and its DATA 111--115 do not
appear, because $G$ never passed through them. The production branch
$R$ may run ahead of $G$; the global chain records only what was
accepted.

\begin{definition}[\defterm{Global finality}]
\label{def:global-finality}
A DATA block $B$ is \defterm{globally final at Espresso height $q$} when
$B\in G_q$, the global chain output by the rule of \cref{sec:exports} at
$q$. Observer $u$ treats $B$ as globally final once it has authenticated
the finalized Espresso history through $J_q$ and the complete L1 control
history that applies to it, and has applied the rule through a step that
appends $B$.
\end{definition}

The rule only appends, so $G_{q-1}$ is a prefix of $G_q$ for every $q$
(\cref{prop:global-prefix}). No later Force removes a globally final
block $B$ or changes the record $\sigma_G(B)$ bound to it, and no extra waiting period
follows. Espresso inclusion alone, or a receipt alone, does not make a
block globally final: an old-version block may be finalized in Espresso
as raw bytes and still never enter $G$.

This is finality of the input prefix. Espresso neither executes the
native transactions nor proves their resulting state; the regional
validity contract and the execution-verification interface are
responsible for that. Finality of a state claim or withdrawal on L1 is a
later, separate stage. Publication lag never permits skipping globally
final input.

\subsection{The global finality rule}
\label{sec:exports}
The global finality rule computes the global chain and the selection
information for its next DATA. It is a deterministic function of its
two input streams, so every verifier that reads the same history
computes the same output.

In plain words, for each new finalized Espresso block the rule does two
things, in this order.
\begin{enumerate}[label=(\arabic*),itemsep=1pt]
\item If an L1 update activates at this block and targets the current
branch version, apply the Force: close the version, install the new map,
keep $G$, and select the next session. If no region is eligible, the
next session stays pending.
\item Append, one at a time, each published DATA block that extends
$G$, belongs to the export session and has a valid regional receipt.
Its Espresso reference must lie at or above the session's reference
floor and strictly before the current block. When the export session's
terminal DATA is appended, select the next session by a regular
handoff.
\end{enumerate}
The rest of this subsection states the state, inputs and output, and
then gives the rule precisely as \cref{alg:global-rule}.

Between heights, the verifier retains the minimal \defterm{replay state}
\[
 \mathcal X_q=(\mathrm{Out}_q,S_{\beta_q},\kappa_q),
\]
where $\mathrm{Out}_q$ is the public output defined below, $\beta_q$
is read from its tail $\mathsf{Tail}_q$, $S_{\beta_q}$ is that branch's
immutable complete map, and $\kappa_q$ is the processed control cursor
of \cref{sec:l1}. These are maintained state, not additional
per-height input streams. In particular, no historical sequence of
status maps is kept for successor selection. The public output of the
protocol remains exactly $\mathrm{Out}_q=(G_q,\mathsf{Tail}_q)$.

\paragraph{Input.}
\label{sec:logical-carrier}
For each newly finalized Espresso block $J_q$, Protocol~3 consumes
exactly two logical input streams:
\begin{enumerate}[label=(I\arabic*),itemsep=2pt]
\item \textbf{Espresso-finalized regional results.}  The authenticated
namespace of $J_q$ adds a set $M_q$ of regional objects.  Complete
regional results consist of DATA bodies, FINAL signatures and any
covering path needed to establish regional finality.  The rule does not
depend on how EXPORT records group these objects.  It uses:
 \begin{itemize}[itemsep=1pt]
 \item block bodies
 $B=(\beta,r,k,h,\mathit{type},x,\Hash(B_{\rm parent}),\rho,\xi)$ of type
 DATA or DRAIN, identified by their hash;
 \item FINAL signatures: a signature of a member on the hash of a block
 $D$, in the signed context of $D$'s session.
 \end{itemize}
Write $\mathcal O_q=M_1\cup\cdots\cup M_q$ for everything published
through $J_q$. A DATA block is an append candidate only when the objects
in $\mathcal O_q$ establish $\mathrm{Fin}_q(\nu,B)$ under the current
next-block descriptor.
\item \textbf{The activated update.} From the complete finalized L1
control log, let $\Lambda_q$ contain the valid update whose activation
threshold is first crossed at $J_q$, if one exists. By
\cref{ass:single-activation}, $\Lambda_q$ is empty or a singleton. Its
single update executes if its target is the current branch version and
is otherwise a stale whole-update no-op.
\end{enumerate}
The genesis checkpoint supplies the initial state rather than a third
per-block input: the chain identifier, region order $\mathcal L_R$,
anchor $A_0$, first replay position $q_{\rm init}$, initial map $S_0$
and control prefix. A later trusted checkpoint supplies an earlier
$\mathcal X$ instead. Espresso headers, namespace-completeness
proofs, $\tagg{AUTH}_{\beta,r,0}$ and L1 authentication proofs are validation
dependencies for the two inputs above, not additional logical input
streams.

Thus the per-height transition consumes only $\mathcal O_q$ and
$\Lambda_q$, together with its already persisted $\mathcal X_{q-1}$.
It never rereads the full update log to choose a regular successor.

Proposals, QCs, recovery authorities
$\tagg{AUTH}_{\beta,r,g}$ with $g>0$ and native sidecars may also appear in
$M_q$; the rule ignores them. The outer fields of an EXPORT, such as its
expected base, interval and version, are routing hints and are ignored
too. A malformed or irrelevant object grants nothing and blocks nothing.

The rule uses one test derived from the input. For a complete session
descriptor $\nu$ and a candidate DATA block $B$, write
$\mathrm{Fin}_q(\nu,B)$ when
$\tagg{CheckFINAL}(C_\nu,B,D,\Sigma,\pi)$ returns \emph{true} for
some block $D$ and path $\pi$ whose bodies all lie in $\mathcal O_q$, with
$\Sigma$ the set of FINAL signatures on $D$ in $\mathcal O_q$. The
signatures may come from different EXPORT records and Espresso blocks;
only their total weight counts.

\paragraph{Output.}
For every Espresso height $q\ge q_{\rm init}$ the rule outputs
\[
 \boxed{\mathrm{Out}_q=(G_q,\mathsf{Tail}_q)},\qquad
 \mathsf{Tail}_q\in\{\tagg{NEXT}(\nu),\ \tagg{PENDING}(\omega)\}.
\]
$G_q$ is the global chain with its block-bound records
(\cref{def:global-chain}).  The tail $\mathsf{Tail}_q$ says who is responsible for
the next DATA after $\operatorname{last}(G_q)$:
\begin{itemize}[itemsep=2pt]
\item $\mathsf{Tail}_q=\tagg{NEXT}(\nu)$ means that the fully specified session
$\operatorname{sid}(\nu)$ is the \defterm{export session}: it owns the
next DATA height and its descriptor supplies the region, anchor,
terminal height, Espresso reference floor and cause. Its accepted
count, the number of its DATA already in $G_q$, is
\[
 a(G_q,\nu)=
 \bigl|\{B\in G_q:\operatorname{sid}(\sigma_{G_q}(B))
                 =\operatorname{sid}(\nu)\}\bigr|.
 \]
\item $\mathsf{Tail}_q=\tagg{PENDING}(\omega)$, where
\[
 \omega=(\mathit{cause},\beta,r',i_{\rm src},A,q^{\rm open}),
\]
means that the next session identifier $(\beta,r')$, its anchor
$A=\operatorname{last}(G_q)$ and its opening reference floor are fixed,
but $\tagg{NextLive}(i_{\rm src},S_\beta;\beta,r')=\bot$: the immutable map
$S_\beta$ has no eligible target region.
\end{itemize}
The current branch version is $\beta(\nu)$ in the \tagg{NEXT} case and
the displayed $\beta$ in the \tagg{PENDING} case.  Exactly one tail
variant is present at every completed replay subposition.

\begin{protocol}[H]{alg:global-rule}{Global finality rule: update replay state and output $\mathrm{Out}_q$}
Initialize $\kappa$ to the control cursor in the genesis checkpoint. Use
$S_0$ and $\tagg{AUTH}_{0,0,0}$ to compute
$i_0=\tagg{NextLive}(\bot,S_0;0,0)$.  If $i_0\ne\bot$, set
\[
 \nu_0=(0,0,i_0,A_0,h(A_0)+\TB_0,q_{\rm init},\tagg{GENESIS}),
 \qquad \mathrm{Out}=(A_0,\tagg{NEXT}(\nu_0));
\]
otherwise set
\[
 \mathrm{Out}=(A_0,
 \tagg{PENDING}(\tagg{GENESIS},0,0,\bot,A_0,q_{\rm init})).
\]
For each finalized block $J_q$ in order, start from $\mathcal X_{q-1}$,
ingest its two inputs---the empty-or-singleton activated update
$\Lambda_q$ and the cumulative
Espresso-finalized object set $\mathcal O_q$ (after adding $M_q$)---and
run the two steps below. The public result is $\mathrm{Out}_q$; the
updated map and control cursor complete $\mathcal X_q$. Write
$\operatorname{last}(G)$ for the current tip.
\begin{enumerate}[label=\textbf{\arabic*)},itemsep=3pt]
\item \textbf{Apply an atomic branch update.}
Let $\beta_0$ be the branch version at the start of the prelude. If
$\Lambda_q=\varnothing$, skip the rest of this step. Otherwise, by
\cref{ass:single-activation}, write
\[
 \Lambda_q=\{U_q\},\qquad
 U_q=(\mathit{chain},c_q,\theta_a,\beta_{\rm target},S_{\rm new}).
\]
Advance $\kappa$ past $U_q$. If $\beta_{\rm target}\ne\beta_0$, the
whole update is stale: leave $\mathrm{Out}$ and $S_{\beta_0}$ unchanged
and skip the rest of this step. Otherwise close $\beta_0$, set
$\beta':=\beta_0+1$ and fix $S_{\beta'}:=S_{\rm new}$. Determine the target
index and source region as follows:
  \begin{itemize}[itemsep=1pt]
  \item if $\mathsf{Tail}=\tagg{NEXT}(\nu)$, write
  $\operatorname{sid}(\nu)=(\beta_0,r)$, let $i_{\rm src}=i(\nu)$ and
  $a=a(G,\nu)$, and set $r'=r$ if $a=0$ and $r'=r+1$ otherwise;
  \item if $\mathsf{Tail}=\tagg{PENDING}
  (\mathit{cause},\beta_0,r',i_{\rm src},A,q^{\rm open})$, keep its $r'$ and
  $i_{\rm src}$.
  \end{itemize}
Let $A=\operatorname{last}(G)$ and
$i'=\tagg{NextLive}(i_{\rm src},S_{\beta'};\beta',r')$.  If $i'=\bot$, set
\[
 \mathsf{Tail}:=\tagg{PENDING}
 (\tagg{FORCE}(c_q,q),\beta',r',i_{\rm src},A,q).
\]
Otherwise set
\[
 \nu'=(\beta',r',i',A,h(A)+\TB_{r'},q,
       \tagg{FORCE}(c_q,q)),
 \qquad \mathsf{Tail}:=\tagg{NEXT}(\nu').
\]
This step never changes $G$ or an existing $\sigma_G(B)$.

\item \textbf{Append matching DATA.}
While $\mathsf{Tail}=\tagg{NEXT}(\nu)$ and $\mathcal O_q$ contains a
DATA block $B$ satisfying, with
$\operatorname{sid}(\nu)=(\beta,r)$,
  \begin{enumerate}[label=(\alph*),itemsep=1pt]
  \item $B$ is signed in this chain's domain by session $(\beta,r)$;
  \item $\mathrm{parent}(B)=\Hash(\operatorname{last}(G))$ and
  $h(B)=h(\operatorname{last}(G))+1\le H(\nu)$;
  \item $\mathrm{Fin}_q(\nu,B)$ holds; and
  \item $\xi(B)$ names a finalized block $J_{q'}$ of the log, with its
  authenticated hash, and $q^{\rm open}(\nu)\le q'<q$,
  \end{enumerate}
append $B$ to $G$ and bind $\sigma_G(B):=\nu$.  If several blocks
qualify, take the one with the least hash.  If $h(B)<H(\nu)$, leave
$\mathsf{Tail}=\tagg{NEXT}(\nu)$ and continue.  If $h(B)=H(\nu)$, write
$\operatorname{sid}(\nu)=(\beta,r)$ and set
\[
 q^{\rm open}_{\rm next}=
 \max\{\operatorname{height}(\xi(B)),q^{\rm open}(\nu)+1\}.
\]
Let $i'=\tagg{NextLive}(i(\nu),S_\beta;\beta,r+1)$. If $i'=\bot$, set
\[
 \mathsf{Tail}:=\tagg{PENDING}
 (\tagg{REGULAR},\beta,r+1,i(\nu),B,q^{\rm open}_{\rm next}).
\]
Otherwise set
\[
 \nu'=(\beta,r+1,i',B,h(B)+\TB_{r+1},
       q^{\rm open}_{\rm next},\tagg{REGULAR}),
 \qquad \mathsf{Tail}:=\tagg{NEXT}(\nu').
\]
Then continue looking for another DATA block.
\end{enumerate}
\end{protocol}

\Cref{alg:global-rule} is the public global-finality rule;
the Force handler in Part~IV~\cite{part4} restates its prelude for the regional runtime. Step~1 consumes the control log and changes only the
branch version, its fixed map and the next-block selection. Step~2
appends exactly the regionally finalized DATA that extends the current
tip and matches that selection. The waiting rules are stated below and
in Part~IV~\cite{part4}.

\paragraph{Why each check.}
\begin{itemize}[itemsep=2pt]
\item \textbf{Export session only} (a). Only the session that owes the
next height may extend $G$. A closed version has no export session, so
after a Force its DATA never enters $G$. A regular successor that runs
ahead becomes the export session only after its predecessor's terminal
DATA is accepted.
\item \textbf{Continuity} (b). $G$ has no gap and no fork. DRAIN never
enters $G$, because it is not DATA and does not advance the height.
\item \textbf{Regional finality} (c). By
\cref{prop:final-kernel-soundness}, a correct replica of the session
committed $B$ after fully validating it. The rule needs no proposal or
recovery history.
\item \textbf{Selection reference} (d). A reference cannot justify itself
through the block that carries it. Because $q^{\rm open}\ge q_f$ for a
session opened by a Force at $J_{q_f}$, no new-version DATA is accepted
in $J_{q_f}$ itself.  A regular successor may already have published
output under its deterministic opening floor; that output can follow the
terminal DATA in the same or a later EXPORT once the terminal enters
$G$.
\end{itemize}


For example, immediately before the Force in the running example,
\[
 G=100\cdot101\cdots108,\qquad
 \mathsf{Tail}=\tagg{NEXT}(\nu_A),\qquad
 \operatorname{sid}(\nu_A)=(7,20),\quad i(\nu_A)=\text{A},\quad
 a(G,\nu_A)=8.
\]
Suppose the executing update has sequence number $c_f$ and activation
block $J_{q_f}$. Step~1 leaves $G$ and its records unchanged and
constructs
\[
 \omega_F=(\tagg{FORCE}(c_f,q_f),8,21,\text{A},108,q_f).
\]
Because B is eligible under the newly installed map $S_8$, Step~1 sets
the tail directly to $\tagg{NEXT}(\nu_B)$, where
\[
 \nu_B=(8,21,\text{B},108,108+\TB_{21},q_f,
        \tagg{FORCE}(c_f,q_f)).
\]
B's new DATA $109'$ carries a reference at least $q_f$, so Step~2 can
accept it from $J_{q_f+1}$ on and bind
$\sigma_G(109')=\nu_B$. A's old 109 is still in $\mathcal O_q$ but
fails (a). If no region were eligible under $S_8$, Step~1 would instead
keep $\tagg{PENDING}(\omega_F)$ until another successful atomic update
created a later version with a new map.




\section{Cross-region rotation}
\label{sec:xprotocols}
This section describes what each regional node does across sessions.
Each node runs a Session Manager. It follows the global chain, regional
results and L1 updates, and starts or stops the node's session programs.
At node level, ideas~2, 3 and~5 of
\cref{sec:core-idea} become four rules:
\begin{itemize}[itemsep=2pt]
\item \textbf{Track three things.} The Session Manager follows the
globally finalized chain $G$, the regionally finalized production chain
$R$, and the update log.
\item \textbf{Regular handoff.} Once $R$ reaches the current session's
terminal DATA with a valid receipt, the selected region relays the
ReadyBundle and enters the next session. Its initial leader directly
proposes a child of that terminal DATA, without waiting for it to enter $G$.
\item \textbf{Force handoff.} When a valid Force activates, the selected
region starts a fresh session from the tip of $G$ retained at activation:
relay the ReadyBundle, enter, and directly propose the next DATA.
\item \textbf{Force-close.} For every session, the Manager terminates its
program and cancels its timers once it verifies that an activated Force
has closed that session's branch version.
\end{itemize}
Both handoffs use the same entry rule: validate the handoff and the
anchor state, relay a ReadyBundle, enter once, and let the initial leader
propose the first DATA directly. They differ only in the anchor. A
regular handoff continues after the predecessor's terminal DATA. A Force
handoff continues after the last DATA retained in $G$.
\Cref{sec:session-manager} lists what the Manager keeps,
\cref{sec:witness} defines both handoffs and the Manager rules, and
\cref{sec:ready} gives the ReadyBundle and the startup protocol.

\paragraph{Example: overlapping handoffs.}
\Cref{fig:concurrent-handoffs} shows the running example. A's terminal
receipt for 110 opens B's regular session $(7,21)$ while a Force is
pending. At $J_{501}$, the Force targets version 7. $G$ is still 108, and
the export session is still A's $(7,20)$. Selecting B again opens a fresh
session $(8,21)$ from 108. B's old $(7,21)$ program stops when B
verifies the Force.

\begin{figure}[tbp]\centering
% Section 6 overview: a regular successor starts while a Force is pending.
% G is held at 108 in this example. Boxes show DATA lineage and local programs,
% not a simultaneous wall-clock transition at all replicas.
\begin{tikzpicture}[
  x=1cm,y=-1cm,font=\fontsize{9}{10.5}\selectfont,
  card/.style={draw,rounded corners=2pt,align=center,
    inner sep=4pt,minimum height=1.05cm},
  kept/.style={card,draw=rrTeal,fill=rrTeal!5,text width=3.1cm},
  old/.style={card,draw=rrBlue!70!black,fill=rrBlue!4},
  fresh/.style={card,draw=rrTeal,fill=rrTeal!5,
    text width=9.7cm,minimum height=1.35cm},
  control/.style={draw=black!35,rounded corners=2pt,fill=black!2,
    align=center,text width=14.6cm,inner sep=5pt},
  force/.style={draw=rrRed,rounded corners=2pt,fill=rrRed!4,
    align=center,text width=9.7cm,inner sep=5pt},
  msg/.style={-{Stealth[length=2mm]},semithick,draw=rrBlue},
  restart/.style={-{Stealth[length=2.2mm]},thick,draw=rrTeal},
  tag/.style={fill=white,inner sep=2pt,align=center,
    font=\fontsize{8.5}{10}\selectfont},
  stop/.style={text=rrRed,align=center,text width=10cm,
    font=\fontsize{9}{10.5}\selectfont}
]
  \path[use as bounding box] (0,-.15) rectangle (15.4,8.25);

  \node[control] at (7.7,.40)
    {An L1 update targets version 7, but has not activated yet.\\
     A's terminal receipt can still start B's regular session.};

  \draw[draw=rrRed!65,densely dashed,rounded corners=3pt]
    (4.55,1.40) rectangle (14.85,3.25);
  \node[tag,text=rrRed] at (9.70,1.40)
    {Old branch suffix ($\beta=7$)};
  \node[kept] (g) at (1.95,2.35)
    {$G:\ 100\to\cdots\to108$\\
     \textbf{Kept by the Force}};
  \node[old,text width=2.5cm] (a) at (6.20,2.35)
    {A: $(7,20)$\\
     $109\to T=110$};
  \node[old,text width=4.0cm] (b) at (12.40,2.35)
    {B: regular $(7,21)$\\
     anchor 110\\
     DATA $111,\ldots,115$};
  \draw[msg] (g.east)--(a.west);
  \draw[msg] (a.east)--node[tag,above] {$\FC_{7,20}(110)$}
    node[tag,below] {regular handoff}(b.west);
  \node[tag,text=rrRed] at (9.70,3.25)
    {Old DATA 109--115 cannot enter $G$ after the Force};

  \node[force] (f) at (9.70,4.25)
    {\textbf{At $J_{501}$: Force activates; $G$ still ends at 108.}\\
     Close version 7; install $S_8$ (A OFF, B/C/D ON); select B.};
  \node[stop] (stopped) at (9.70,5.30)
    {When B verifies the Force: stop $(7,21)$ and cancel its timers.\\
     Stop every other local version-7 program too.};
  \draw[-{Stealth[length=1.5mm]},draw=rrRed] (f.south)--(stopped.north);

  \node[fresh] (newb) at (9.70,6.80)
    {\textbf{New B: $(8,21)$, anchor 108, state $\mathit{st}(108)$}\\[2pt]
     ReadyBundle $\longrightarrow$ Enter $\longrightarrow$
     propose $109',110',\ldots$\\[2pt]
     Fresh locks, votes and timers};
  \draw[restart] (g.south)--(1.95,6.80)--(newb.west);
  \node[tag,text=rrTeal,text width=3.0cm] at (1.95,4.50)
    {Force continues\\from \textbf{108}, not 110};
  \draw[-{Stealth[length=1.5mm]},draw=black!50]
    (stopped.south)--(newb.north);

  \node[align=center,text width=15cm,font=\fontsize{9}{10.5}\selectfont]
    at (7.70,8.00)
    {\textbf{Same region, new session.}\quad
     Regular handoff follows $R$; Force restarts from the retained $G$.};
\end{tikzpicture}

\caption{Regular handoff inherits terminal DATA; Force inherits the global
tip at activation. If a node verifies the Force first, a late terminal
receipt cannot start $(7,21)$. The new session's proposals still need
Protocol~3 acceptance to enter $G$.}
\label{fig:concurrent-handoffs}
\end{figure}

\subsection{What the Session Manager keeps}
\label{sec:session-manager}
\label{sec:runtime-modules}
The Manager tracks three kinds of progress:
\begin{itemize}
\item \textbf{Global chain $G$.} Run \cref{alg:global-rule} on the
verified public inputs and keep its output $(G,\mathsf{Tail})$.
This is the accepted DATA prefix and the selection for its next block.
\item \textbf{Regional chain $R$.} Keep the complete, validated regional
results of the current production branch, grouped by session and linked
by their handoffs. Direct results can extend $R$ before they enter $G$.
Keep older results as history when a Force replaces that branch.
\item \textbf{Updates.} Follow the authenticated L1 control log and its
Espresso activation positions. Keep the current version, its fixed
$S_\beta$, the processed control cursor and pending updates. A regular
handoff uses its version's fixed map; it does not rescan the update log.
\end{itemize}
Authorities and source proofs are validation inputs for this work.
The Manager applies the public rules itself; it makes no independent
choice of region, cutoff or global DATA.

\paragraph{State for the next anchor.}
Write $\mathit{st}(B)$ for the executable state after the exact block
$B$. Each operational node maintains a \defterm{baseline state} at
$\operatorname{last}(G)$. A region selected to take over by regular
handoff also follows the predecessor's regionally finalized DATA in a
\defterm{successor state}. The baseline prepares for a Force; the
successor state prepares for a regular handoff.

For example, let A finish at 110 while $G$ is only 108. B prepares
$\mathit{st}(110)$ for regular handoff and keeps $\mathit{st}(108)$ as
its baseline. A Force that retains 108 uses the latter. If $G$ advances
before the Force, the retained anchor advances with it. Once a handoff
fixes an anchor, the program takes its own matching state; subsequent
baseline work cannot change that anchor. A node whose execution lags
must finish preparing the exact state before entering.

\paragraph{Session programs.}
\label{sec:session-lifecycle}
A program runs the regional engine for one pair $(\beta,r)$. It keeps
that session's view, locks, votes, QCs, timers and executable state.
Another pair has independent core state, even if it reuses $r$.
The Manager supplies verified references and authorities and receives
the program's results and completion event.

A program is \defterm{Assigned} before initial Enter,
\defterm{Running} while its engine is Active or Frozen, and
\defterm{Terminated} after Completion or Force-close. Quit and recovery
stay inside Running. The Manager durably records the fixed handoff,
whether initial Enter occurred, its time, and any termination. It also
keeps version-wide closed marks. These local records prevent repeated
entry or recreation after termination; they are not the block-bound
records $\sigma_G(B)$ in $G$.

\subsection{Handoffs and session management}
\label{sec:witness}
A \defterm{handoff} fixes why a new session starts, who runs it and
which DATA it continues after. Its fields are the descriptor $\nu$ of
\cref{sec:global-chain}, the target authority $\tagg{AUTH}_{\beta,r,0}$,
the version's map $S_\beta$, the L1 reference floor and the evidence
that establishes them. The anchor keeps its original producing pair,
hash and commitment evidence. A node never replaces an installed
handoff with different fields for the same pair.

\paragraph{Regular handoff.}
A complete valid receipt $\FC_{\beta,r}(T_{\beta,r})$ fixes the
successor's anchor at $T_{\beta,r}$ and its pair at $(\beta,r+1)$.
The selected region and Espresso reference floor are
\[
 i'=\tagg{NextLive}(i_{\beta,r},S_\beta;\beta,r+1),\qquad
 q^{\rm open}_{\beta,r+1}
 =\max\{\operatorname{height}(\xi(T_{\beta,r})),
         q^{\rm open}_{\beta,r}+1\}.
\]
NextLive is the selector of \cref{sec:l1}. This handoff uses the
unchanged $S_\beta$, regardless of when a node receives the receipt.
It can create and start the successor before the terminal DATA enters
$G$. A terminal QC alone is insufficient. If the reference floor names
a header not yet available, entry can still occur; the first DATA waits
for a valid finalized reference at or above that floor.

\paragraph{Force handoff.}
\label{sec:retry-slot}
When an update executes at $J_{q_f}$, \cref{alg:global-rule} closes the
old version $\beta$, installs $S_{\beta+1}$ and retains the preceding
global prefix $G^\star$. Its new anchor is
$A=\operatorname{last}(G^\star)$, with $q^{\rm open}=q_f$.
If the old tail is $\tagg{NEXT}(\nu)$, let $(\beta,r)$ be its export
session and let $a=a(G^\star,\nu)$. The target index is
\[
 r'=\begin{cases}
 r,&a=0,\\
 r+1,&a>0.
 \end{cases}
\]
The source region is $i(\nu)$. If the old tail is PENDING, keep its
stored source region and target index. Select under the new map and
$\tagg{AUTH}_{\beta+1,r',0}$. The Manager uses the resulting pair
$(\beta+1,r')$ and exact anchor; it does not choose from the furthest
speculative program or the $G$ it happens to see later.

A PENDING result opens no program. Missing authority means wait for
evidence, not that all regions are ineligible. A later successful update
may replace a pending version, but passage of time alone cannot do so.
A stale update changes neither the map nor any program.
Genesis uses the initial checkpoint's selection and anchor.

\paragraph{Verify the handoff, then check permission to enter.}
\label{sec:verify-handoff}
The procedure $\tagg{VerifyHandoffHistory}$ authenticates the fixed
fields from a trusted replay checkpoint and the required source
history. For REGULAR it checks the predecessor's boundary, terminal
height and complete terminal receipt, then recomputes the target and
reference floor. For Force it checks the authorized update, its unique
activation block, the matching target version and the exact retained
$G^\star$ under \cref{alg:global-rule}. It verifies every applicable
control needed for that history, not just a sender-selected update.
An inherited receipt may use \tagg{CheckFINAL}; the target's own
authority, state and admission still require their full checks.

For either handoff, authenticate $\tagg{AUTH}_{\beta,r,0}$ for the exact
target pair, its fixed configuration, quota and map,
the anchor's exact DATA and original commitment evidence, and
\[
 \lambda^{\rm open}_{\beta,r}
 =\max\{\lambda(A_{\beta,r,0,0}),
         \lambda(\tagg{AUTH}_{\beta,r,0})\}.
\]
The authority term is its L1 source-height bound, including required
publication, activation and branch-control heights (\cref{sec:authority-context}).
Incomplete evidence causes a wait. Invalid or conflicting fixed fields
are rejected. Historical validation can still verify retained DATA
from a closed version; it grants no right to restart its producer.

\label{sec:can-enter}
The separate check $\tagg{CanEnter}$ requires that the target version
is current and not marked closed in the node's verified history,
the fixed handoff selects the node's region, and the pair has neither
entered nor terminated. Both checks apply to every creation and entry
path, including a ReadyBundle received before a local program exists.

\begin{protocol}[H]{alg:rotation}{Session Manager rules}
Each correct regional node handles one event at a time. Every action
below uses authenticated inputs and its persisted local records.
\begin{enumerate}[label=\textbf{\arabic*)},itemsep=3pt]
\item \textbf{Track.} Process public inputs under
\cref{alg:global-rule}, update $G$ and the control state, and validate
regional results for $R$. Advance the baseline and, when this region is
the regular successor, its successor state. This work signs no BFT
message.
\item \textbf{Create.} If a verified regular handoff or the public
Force/genesis selection assigns the node's region, check
\tagg{CanEnter} and record the fixed handoff in an Assigned program.
A valid ReadyBundle can supply the same assignment. Repeated inputs
reuse the record; a different handoff for that pair is rejected.
\item \textbf{Start.} When the Assigned program is Ready, run
\cref{alg:bootstrap} once. Buffered messages cannot authorize an earlier
vote, and a late callback cannot bypass the entry checks.
\item \textbf{Force-close.} On first verifying that an effective Force closed
version $\beta$, atomically persist its closed mark and terminate every
local program of that version. Cancel its BFT timers, progress checks,
proposals, votes, term changes and recovery. Assigned programs are
cancelled too. Every later callback and BFT action checks the closed
mark, including for a pair whose program was never created. Keep $G$
and its records unchanged; reset the current production branch to the
Force's retained $G^\star$. No timer-drain period or close-bundle
exchange is required.
\item \textbf{Complete.} When a Running program performs Completion in
\cref{alg:steady}, record normal termination. Its complete terminal
receipt is relayed and its result offered for publication by that
existing rule; no further closing protocol runs.
\item \textbf{Serve and persist.} Supply Running programs with verified
references and authorities. Offer complete results to the Publish Layer
and send terminal results to the selected regular successor. Keep
required data, evidence and outstanding sends independently of program
lifetime, under \cref{app:storage}. Persist entry and termination before
dispatching signatures; duplicates and restarts repeat no logical action.
\end{enumerate}
\end{protocol}

\paragraph{What Force-close means.}
\label{sec:force-close}
The public cutoff is fixed by \cref{alg:global-rule}; a local program
stops when its Manager verifies that cutoff. Nodes need not learn it
at the same real time. The closed mark immediately disables all old
BFT actions, even if physical resource cleanup takes longer. New
programs need not wait for that cleanup.

The protocol preserves all of $G^\star$. Old-version DATA outside it
can never be appended later, including DATA that was regionally final
before the Force. No regional agreement or completion guarantee is
claimed for this discarded suffix during forced termination. An old
node that has not learned the Force may briefly continue, but its late
output cannot cross the public cutoff or change a new program's core
state. Already formed evidence remains historical evidence; aborting a
program creates no FINAL signature and supplies no terminal receipt.

If a terminal receipt arrives first, Completion may start a regular
successor in the old version. A later verified Force terminates that
successor too. If the Force is processed first, the old receipt is kept
as evidence and cannot restart an old program. Termination does not
delete retained DATA, authorize pruning or release stake.

\Needspace{15\baselineskip}
\subsection{ReadyBundle and direct startup}
\label{sec:ready}
\begin{definition}[\defterm{Ready}]
\label{def:ready}
An Assigned program for $(\beta,r)$ is Ready when:
\begin{enumerate}[label=(\roman*),itemsep=2pt]
\item its handoff passes \tagg{VerifyHandoffHistory} and
\tagg{CanEnter};
\item the exact anchor, its original producing identity, full DATA,
commitment evidence and required ancestry have been validated;
\item the node holds a validated executable state after that exact
anchor, with no remaining reconstruction or installation work; and
\item its initial authority, fixed membership, quota, map and reference
floors, with their validation dependencies, are available.
\end{enumerate}
\end{definition}
The Espresso header at the opening floor need not exist yet.
\Cref{sec:blocks} still requires a suitable finalized header before the
first DATA proposal. State prepared for a different anchor never
satisfies Ready.

\paragraph{The ReadyBundle.}
\label{sec:ready-bundle}
The ReadyBundle replaces a startup voting round. The first correct
member to enter sends everything the others need to enter. So all
correct members enter within the regional delay bound
$\Delta_{i_{\beta,r}}$ of the first one
(\cref{prop:sync-start}), and the first leader can propose at once.

For a target $(\beta,r)$, a ReadyBundle is an authenticated message to
its fixed regional membership. It contains the fixed handoff and all
the material in \cref{def:ready}: exact anchor DATA and evidence,
the already constructed executable state $\mathit{st}(A)$ with its
native validation evidence, the authority and bounds, and the complete
required source and validation dependencies. It carries no inherited
locks, votes or timers.

A sender supplies the full bytes; a state root, URL or promise to fetch
later is not a replacement. Recipients verify the material themselves.
They may reuse a local state only after verifying that it is the state
after the same exact anchor. The received state must be directly
usable; a snapshot that still needs unbounded import does not satisfy
the complete-message contract below.

At its first Enter, every correct member sends its complete ReadyBundle
to all members. Sending is part of that entry event and is not deferred
to a later task. A recipient validates the handoff and public context,
checks its current closed marks, and, if it has not entered or
terminated, adopts the exact state and enters in the same bounded
processing event. It sends its own complete bundle at that event.
Duplicates never reset entry time or core state. Closed or terminated
targets are not entered. There is no READY quorum or new activation
certificate.

\begin{protocol}[H]{alg:bootstrap}{Enter and propose the first DATA}
A member $v$ of the selected region runs these steps for one fixed
session $(\beta,r)$.
\begin{enumerate}[label=\textbf{\arabic*)},itemsep=3pt]
\item \textbf{Enter and relay.} Once Ready, atomically recheck
\tagg{CanEnter}, set the durable initial-entry flag, install the fixed
anchor, authority, quota, reference floors and exact executable state,
and set the entry lock to $Q^0_{\beta,r}$. Enter $(\beta,r,0,0)$
unfrozen, mark the program Running and record $e_v(\beta,r,0,0)$.
Send the complete ReadyBundle in this event, start the progress checks
of \cref{alg:recovery}, and process buffered evidence before any vote
or proposal.
\item \textbf{Propose directly.} If $v=L(\beta,r,0,0)$, the program is
Running and view $(\beta,r,0,0)$ is still current and unfrozen,
run Propose in \cref{alg:steady} as soon
as its required references and other proposal inputs are available.
At $(g,t,j)=(0,0,0)$ use
\[
 B_{\rm parent}=A_{\beta,r,0,0},\qquad
 Q_{\rm parent}=Q^0_{\beta,r}.
\]
The first DATA has height $h(A_{\beta,r,0,0})+1$ and physical position
1. No initial STATUS, NEWVIEW, ANCHOR vote, ANCHOR QC or bootstrap
waiting period precedes it. Members vote only after Enter, using the
normal full-validation and commit-timer rules.
\end{enumerate}
\end{protocol}
A first proposal carries the handoff and base justification needed to
validate its fixed opening. Later DATA use the normal parent QCs.
Recovery within the same session ($g>0$) still uses STATUS, NEWVIEW,
the prescribed waits and an ANCHOR QC: it must account for that
session's existing locks. The progress-check deadlines are unchanged.

The synchronous-start guarantee is stated and proved in
\cref{prop:sync-start}. Let $e$ be the first correct initial Enter for
$(\beta,r)$ and write $\Delta=\Delta_{i_{\beta,r}}$. Under the
proposition's premises, every correct selected member enters by
$e+\Delta$.
This bound uses no first-member election or common startup deadline.
If a correct member votes at time $s$, then $s\ge e$; hence the other
correct members can enter and process its complete proposal by
$s+\Delta$, subject to the same applicability and processing premises.
This is the entry condition used by Part~I's vote-window argument.
Native work outside that processing bound requires its own time budget;
calling it installation does not make it part of $\Delta$.

\paragraph{Preparing the first bundle.}
\begin{assumption}[Bounded preparation for a selected live region]
\label{ass:ready-service}
Let $J_o$ be the block in which public replay first assigns an export
session to a region, and let $t_0=t_E(o)$. If the selected correct
members stay online, the declared input, load and state-service
conditions hold, and no Force supersedes the handoff, then by
$t_0+T_{\rm ready}$ each member has either entered or completed all
preparation and validation needed for its guarded Enter and complete
ReadyBundle.
\end{assumption}
For a regular successor, $J_o$ is the block where the predecessor's
terminal DATA enters $G$; direct delivery may have started the program
earlier. For a Force successor, $J_o$ is the activation block.
A PENDING boundary has no selected region and starts no preparation
obligation. The budget includes source observation, retrieval and native
preparation from $t_0$; it does not restart when a node first downloads
the evidence. It neither supplies a bound on arbitrary offline catch-up
nor creates a new Force trigger.

\section{Examples and progress conditions}
\label{sec:examples}
The first six subsections walk through examples: a regular rotation, a
Force, concrete view changes, a session with no output, repeated Forces,
and the exact cutoff. The final two subsections state the conditions
needed for progress (\cref{sec:progress,sec:global-progress}).
The security and lifecycle propositions and their proofs are collected
in \cref{app:security}.

\subsection{Regular rotation and a late DRAIN}
\label{sec:regular-example}
Let $h(A_{\beta,r,0,0})=100$, $\TB_r=10$ and $K=4$. The DATA terms produce
101--104, 105--108 and 109--110. A terminal QC below 110 moves to the next
term. A receipt for DATA 110 opens the next region, which starts DATA 111
after Ready and Enter. Step by step:
\begin{enumerate}
\item A commits 110 under its regional BFT rules, if necessary through a
DRAIN during recovery.
\item A's members collect the complete regional receipt for 110 and send
it to B.
\item B verifies that it is selected and obtains the data of 110, the
execution state after 110 and the target authority.
\item B's members relay ReadyBundles and enter $(7,21,0,0)$.
Its initial leader proposes 111 directly, using the fixed anchor's
base justification.
\item A's result is published to Espresso in parallel; $G$ may reach 110
later.
\end{enumerate}
Step 2 needs the receipt. A QC for 110 alone is not enough for B to take
over.

\begin{figure}[htbp]\centering
% Regular handoff: a terminal receipt transfers the turn from A to B.
% Representative arrows only; time runs downward and is not to scale.
\begin{tikzpicture}[
  x=1cm,y=-1cm,font=\fontsize{9}{10.5}\selectfont,
  party/.style={draw,rounded corners=2pt,fill=white,align=center,
    inner sep=3pt,minimum width=2.2cm},
  lane/.style={draw=black!30,densely dashed},
  msg/.style={-{Stealth[length=1.7mm]},semithick,draw=rrBlue},
  direct/.style={msg,densely dashed},
  tag/.style={fill=white,inner sep=1.5pt,font=\fontsize{8}{9}\selectfont},
  note/.style={anchor=north west,align=flush left,text width=5.35cm,
    inner sep=0pt,font=\fontsize{8.5}{10}\selectfont},
  phase/.style={anchor=west,fill=white,inner sep=2pt,
    font=\fontsize{9}{10.5}\selectfont\bfseries},
  dot/.style={circle,fill=black,inner sep=1.1pt},
  timer/.style={line width=2.6pt,draw=rrTeal!75},
  wait/.style={line width=2.6pt,draw=rrBlue!70!black,
    dash pattern=on 2.5pt off 1.5pt},
  entry/.style={draw=rrBlue,line width=1.2pt,dotted}
]
  \def\xA{1.2}\def\xR{3.7}\def\xL{6.2}\def\xB{8.7}\def\xn{9.65}
  \path[use as bounding box] (-.1,-.65) rectangle (15.05,13.05);
  \node[party] at (\xA,0) {A: leader /\\receipt holder};
  \node[party] at (\xR,0) {A: other\\replicas};
  \node[party] at (\xL,0) {B: initial leader\\session $(\beta,r+1)$};
  \node[party] at (\xB,0) {B: other\\replicas};
  \foreach \x in {\xA,\xR,\xL,\xB}
    \draw[lane] (\x,.55)--(\x,12.95);

  \node[phase] at (0,.85) {1.\ A commits terminal DATA in session $(\beta,r)$};
  \draw[msg] (\xA,1.40)--node[tag,above,sloped] {$T_{\beta,r}=110$}(\xR,1.70);
  \draw[msg] (\xR,2.15)--node[tag,above,sloped] {DATA votes}(\xA,2.45);
  \foreach \x/\s/\e in {\xA/1.40/3.80,\xR/1.70/4.10} {
    \draw[timer] (\x+.20,\s)--(\x+.20,\e);
    \node[dot] at (\x,\s) {};
    \draw[thick] (\x-.12,\e)--(\x+.45,\e);
  }
  \node[note] at (\xn,1.35) {Replicas validate and vote for the last DATA
    block. Each starts its own $2\Delta_{\rm A}$ commit timer.};
  \node[tag,anchor=east] at (\xA-.15,2.90) {QC$(110)$};
  \node[note] at (\xn,2.80) {The terminal QC does not hand the turn to B.
    The quota permits no DATA 111 in A's session.};
  \draw[msg] (\xR,4.25)--node[tag,above,sloped] {FINAL$(110)$}(\xA,4.55);
  \node[dot] at (\xA,4.55) {};
  \node[note] at (\xn,3.95) {Timers expire; replicas commit 110 and send
    FINAL. Matching signatures of weight $>F_{\rm A}$ form
    $\FC_{\beta,r}(T_{\beta,r})$. Receipt holders complete the old engine.};

  \node[phase] at (0,5.55) {2.\ Direct transfer of the regular boundary};
  \node[note] at (\xn,5.50) {Verify the session and quota. Select B from
    the fixed map $S_\beta$ under $\tagg{AUTH}_{\beta,r+1,0}$. The handoff keeps
    $\beta$ and changes $r$ to $r+1$.};
  \draw[direct] (\xA,6.35)--node[tag,pos=.55,above,sloped]
    {REGULAR witness + bundle}(\xL,7.35);
  \draw[direct] (\xA,6.35)--(\xB,7.55);
  \node[note] at (\xn,6.90) {Direct transfer may precede Espresso
    publication. It waits for neither global acceptance nor a future
    reference header.};

  \node[phase] at (0,8.05) {3.\ B enters and proposes directly};
  \node[dot] at (\xL,8.45) {};
  \draw[msg] (\xL,8.45)--node[tag,above,sloped]
    {ReadyBundle}(\xB,8.90);
  \node[dot] at (\xB,8.90) {};
  \draw[msg] (\xB,9.15)--node[tag,above,sloped]
    {ReadyBundle relay}(\xL,9.40);
  \node[note] at (\xn,8.40) {Verify the fixed handoff, exact state after
    110, authority and bounds. Enter once and relay the complete bundle.
    No initial STATUS or ANCHOR exchange.};
  \draw[msg] (\xL,10.00)--node[tag,above,sloped]
    {DATA 111 + $Q^0_{\beta,r+1}$}(\xB,10.30);
  \foreach \x/\s/\e in {\xL/10.00/12.40,\xB/10.30/12.70} {
    \draw[timer] (\x+.20,\s)--(\x+.20,\e);
    \node[dot] at (\x,\s) {};
    \draw[thick] (\x-.12,\e)--(\x+.45,\e);
  }
  \draw[msg] (\xB,10.85)--node[tag,above,sloped]
    {DATA votes}(\xL,11.15);
  \node[note] at (\xn,10.00) {The first proposal uses
    $Q^0_{\beta,r+1}$ and a finalized reference at or above the opening
    floor.};
  \node[note] at (\xn,11.75) {Normal DATA votes start $2\Delta_{\rm B}$
    commit timers. If uninterrupted, the timers commit 111 and send
    FINAL. Recovery keeps its own ANCHOR QC.};
\end{tikzpicture}

\caption{Regular cross-region rotation: a terminal FINAL receipt fixes
the next session, whose leader proposes directly after entry. The immutable branch map $S_\beta$ fixes the successor, while
the terminal Espresso reference fixes only the next reference floor.
Regional production and Espresso publication proceed in parallel. Neither
the terminal receipt nor the new region waits for global confirmation.}
\label{fig:regular-cross-region}
\end{figure}
\FloatBarrier

Now suppose 110 has a QC but a Quit cancels its timer. Recovery opens at
height 110 and proposes a DRAIN, which supplies a new timer without using
DATA quota. When the DRAIN commits the terminal DATA, replicas send FINAL
for 110, not for the DRAIN. The regular successor still inherits 110. The
old L1 and Espresso references of 110 say nothing about when it was
committed or exported.

Suppose $G$ is only at 100 while A has committed 110 and B has committed
111. The Publish Layer can export 101--111 with the required A and B
receipts and the regular witness. Once those DATA are accepted from
finalized Espresso history, the whole interval is global. A record that
declares only 111 can still supply the predecessors first if it also
carries their DATA and receipts. It cannot skip ancestors that are truly
missing.

\subsection{A multisignature Force and the retained anchor}
\label{sec:force-example}
This is the running example of \cref{sec:overview}. The region order is
A, B, C and D. A's session $(7,20)$ produces DATA 101--110. Its terminal receipt
lets B start $(7,21)$ before all of A's DATA is globally accepted. B has
produced through 115, while $G$ ends at 108. The export session is still
A's: the next required DATA is 109.

The committee submits one atomic update with activation threshold 120
that satisfies advance announcement. It targets branch version 7 and
supplies a complete new map that sets A OFF and keeps B, C and D ON.
Before the first finalized Espresso block that
reaches 120, valid old output may still extend $G$.

If $G$ still ends at 108 at that block, replay checks the target and
atomically closes version 7 and installs the new map with A OFF before
the payload. It keeps 108. Old
A:109--110 and B:111--115 can no longer be appended. Their receipts stay
as historical evidence of regional commitment.

\begin{center}\small
\begin{verbatim}
Production before the Force:
  A: 101 ... 108 109 110
  B:                     111 ... 115
Kept in G at the Force:
     101 ... 108
New B after the Force:
                new 109, new 110, ...
\end{verbatim}
\end{center}

If B is eligible, the forced target is $(8,21)$, because A has already
contributed accepted DATA. The new engine enters $(8,21,0,0)$ from 108 and
may propose DATA 109 directly after Ready and Enter, using $Q^0_{8,21}$. It does not
wait for Espresso results to be registered on L1. New-version DATA is
accepted only in a later Espresso block.

If an older update had already moved version 7 to 8 before 120, this
update would be stale in full: it would neither close version 8 nor
install its map. If instead A had completed its export and B had become
responsible through regular rotation, the version would still be 7: the
instruction would execute, keep the newer $G$, and select after the
actual export region B.

With quota ten, old B $(7,21)$ starts after 110 and ends at 120. New B
$(8,21)$ starts after 108 and ends at 118. They share turn index 21 and
its quota, but have distinct authority contexts, engines and anchors;
the selected snapshots and leader-rule inputs may differ.

\begin{figure}[p]\centering
% Atomic update, immediate old-program abort and direct initial proposal.
% beta is the numeric branch version, and (beta,r) identifies a session.
% Vertical spacing illustrates event order, not elapsed time.
\begin{tikzpicture}[
  x=1cm,y=-1cm,font=\fontsize{8.3}{9.8}\selectfont,
  party/.style={draw,rounded corners=2pt,fill=white,align=center,
    inner sep=3pt,text width=1.22cm,minimum height=.66cm},
  lane/.style={draw=black!28,densely dashed},
  msg/.style={-{Stealth[length=1.6mm]},semithick,draw=rrBlue},
  read/.style={msg,densely dashed},
  force/.style={msg,draw=rrRed},
  tag/.style={fill=white,inner sep=1.5pt,font=\fontsize{7.7}{9}\selectfont},
  note/.style={anchor=north west,align=flush left,text width=5.55cm,
    inner sep=0pt,font=\fontsize{8.3}{9.8}\selectfont},
  phase/.style={anchor=west,fill=white,inner sep=2pt,
    font=\fontsize{9}{10.5}\selectfont\bfseries},
  dot/.style={circle,fill=black,inner sep=1.1pt},
  wait/.style={line width=2.5pt,draw=rrTeal!75,
    dash pattern=on 2.5pt off 1.5pt},
  rulebox/.style={draw=rrRed!55,rounded corners=2pt,fill=rrRed!3,
    align=center,text width=7.9cm,inner sep=5pt}
]
  \def\xC{.75}\def\xI{2.3}\def\xE{3.85}
  \def\xN{5.4}\def\xL{6.95}\def\xB{8.5}\def\xn{9.45}
  \path[use as bounding box] (-.05,-.55) rectangle (15.05,16.05);
  \node[party] at (\xC,0) {Multisig};
  \node[party] at (\xI,0) {L1};
  \node[party] at (\xE,0) {Espresso};
  \node[party] at (\xN,0) {Node /\\certifier};
  \node[party] at (\xL,0) {B: initial\\leader};
  \node[party] at (\xB,0) {B: other\\replicas};
  \foreach \x in {\xC,\xI,\xE,\xN,\xL,\xB}
    \draw[lane] (\x,.52)--(\x,15.65);

  \node[phase] at (0,.88) {1.\ Announce an atomic branch update on L1};
  \draw[force] (\xC,1.55)--(\xI,1.85);
  \node[tag,anchor=south] at (1.52,1.48) {submit $U$ at $s$};
  \node[note] at (\xn,1.30) {$U$ names target version $\beta$, future
    time $\theta_a$, and a complete map $S_{\rm new}$. If $\beta$ is
    stale, the entire update is a no-op.};
  \node[dot] at (\xI,2.65) {};
  \draw[read,draw=rrRed] (\xI,2.65)--node[tag,above,sloped]
    {finalized update + complete history}(\xN,3.30);
  \node[dot] at (\xN,3.30) {};
  \node[note] at (\xn,2.85) {By $s+T_{\rm L1}$, all correct participants
    can verify the finalized update. $T_{\rm L1}$ includes dissemination
    and verification. This precedes the Espresso crossing.};
  \node[tag,align=center] at (5.35,4.15)
    {Before activation: version $\beta$ and $S_\beta$ stay fixed;\\
     normal regional handoffs may continue.};
  \node[phase] at (0,5.25) {2.\ Finalize the Espresso boundary; apply one update};
  \draw[force] (\xE,5.95)--(\xN,6.25);
  \node[tag,anchor=south] at (4.62,5.90) {$J_{q^*}$ finalized};
  \node[note] at (\xn,5.75) {$q^*=\min\{q:\theta_q\ge\theta_a\}$ for
    this update. Interpret the complete history through $J_{q^*-1}$
    first.};
  \node[rulebox] at (4.65,7.30) {At the start of $J_{q^*}$:\\[2pt]
    $\Lambda_{q^*}$ is empty or contains one update;\\
    matching target: close $\beta$ and install full $S_{\rm new}$;\\
    stale target: the whole update is a no-op;\\
    only then process EXPORT records.};
  \node[note] at (\xn,7.05) {Health detection belongs to multisig policy.
    A report or missing DATA cannot execute Force. No Ready or output
    condition gates a valid matching update.};
  \node[dot,fill=rrRed] at (\xN,8.55) {};
  \node[tag,anchor=east,text=rrRed] at (\xN-.12,8.55)
    {anchor $:=\operatorname{last}(G^\star)$};
  \node[note] at (\xn,8.70) {Close version $\beta$ for appends; keep $G$.
    Advance to $\beta+1$. Use $S_{\beta+1}$ and NextLive to select B,
    or keep a PENDING tail.};
  \draw[read] (\xN,9.60)--node[tag,pos=.55,above,sloped]
    {complete boundary}(\xL,10.00);
  \draw[read] (\xN,9.60)--(\xB,10.15);
  \node[note] at (\xn,10.15) {If the old export session has accepted
    count $a$: keep $r$ when $a=0$, otherwise use $r+1$.
    New-version DATA can append only in a later Espresso block.};

  \node[phase] at (0,11.10) {3.\ B enters the new session and proposes directly};
  \node[dot] at (\xL,11.65) {};
  \draw[msg] (\xL,11.65)--node[tag,above,sloped]
    {ReadyBundle}(\xB,12.00);
  \node[dot] at (\xB,12.00) {};
  \draw[msg] (\xB,12.35)--node[tag,above,sloped]
    {bundle relay}(\xL,12.60);
  \node[note] at (\xn,11.65) {B verifies the handoff, exact anchor state
    and $\tagg{AUTH}_{\beta+1,r',0}$, then enters $(\beta+1,r',0,0)$ and relays the
    complete bundle.};
  \draw[msg] (\xL,13.20)--node[tag,above,sloped]
    {first DATA}(\xB,13.55);
  \draw[msg] (\xB,14.10)--node[tag,above,sloped]
    {DATA votes}(\xL,14.35);
  \node[note] at (\xn,13.20) {The first DATA extends
    $A=\operatorname{last}(G^\star)$ using $Q^0_{\beta+1,r'}$.
    No initial ANCHOR QC is needed. Normal DATA votes and commit timers
    follow.};
  \node[note] at (\xn,15.00) {Each Manager aborts its old-version
    programs and timers when it verifies the Force. The new session
    does not wait for an old timer drain.};
\end{tikzpicture}

\caption{Explicit Force: L1 supplies the instruction, Espresso time fixes
its cutoff, and the next region inherits exactly the $G$ at the start of
the activating block. There is no automatic trigger and no second
Espresso-to-L1 wait.}
\label{fig:forced-cross-region}
\end{figure}
\FloatBarrier

\subsection{Concrete view updates}
\label{sec:force-view-examples}
The following examples assume complete authority and an eligible target.
Every Force has an explicit valid MSIG cause.

\paragraph{A has accepted output.}
A is in $(7,20,2,3)$, producing 101--110, while $G=108$. Its export
session is 20 with a positive accepted count. Selecting B gives
\[
 (7,20,2,3)\xrightarrow{\mathrm{Force}}(8,21,0,0).
\]
B continues after 108. Only the new engine's $g$ and $t$ start at zero;
the retained DATA keeps its original provenance.

\paragraph{Regular rotation has already reached B.}
A first hands off
\[
 (7,20,0,2)\xrightarrow{\mathrm{regular}}(7,21,0,0).
\]
B reaches $(7,21,0,1)$ and contributes DATA through 113 to $G$. The next
DATA, 114, belongs to B's session 21. A Force selecting C gives
\[
 (7,21,0,1)\xrightarrow{\mathrm{Force}}(8,22,0,0).
\]
C continues with DATA 114. A scheduled MSIG targeting version 7 still
matches, because regular rotation did not change $\beta$. Even if the
committee saw only $G=108$ when it signed, it cannot make the chain go
back to 108.

\paragraph{Speculative production is ahead of the export session.}
A's $(7,20)$ has regional DATA 101--110; B's $(7,21)$ has 111--120; C runs
$(7,22,1,2)$. If $G=108$, the next required DATA is still A's 109. A Force
selecting B opens $(8,21,0,0)$, not $(8,23,0,0)$. It closes the entire old
version, including C's speculative program. Using a lower turn index than
the latest speculative one removes no globally accepted DATA.

\paragraph{A stale instruction changes no view field.}
One MSIG targets version 7 at threshold 120. An older update has already
moved the version to 8 at an earlier position. Suppose the new
engine is now in $(8,21,1,2)$. Processing the stale instruction leaves
\[
 (8,21,1,2)\xrightarrow{\mathrm{stale\ MSIG}}(8,21,1,2).
\]
Its proposed map does not take effect. To change the map or force again,
the committee needs another valid atomic update that targets version 8.

\subsection{Zero output, an offline region and explicit exclusion}
\label{sec:retry-example}
A's $(7,20)$ opens after DATA 100 but gets no DATA accepted. The committee
observes the failure and submits one update that sets A OFF and carries
the complete new map while targeting version 7. At its activation, if
$G$ is still 100, the target index stays 20. NextLive can choose B, which
enters $(8,20,0,0)$ with a fresh engine. Reusing the number 20 does not
mean selecting A again.

If B also adds no DATA, another explicit MSIG can choose C in
$(9,20,0,0)$. Neither B's Ready nor a first DATA is required before that
instruction. If A's DATA 101 had been accepted before the first Force,
the new index would be 21 with anchor 101.

The monitoring report itself has no effect.
\begin{center}\small
\begin{tabularx}{\linewidth}{@{}>{\raggedright\arraybackslash}p{.45\linewidth}Y@{}}\toprule
Situation & What the protocol does\\\midrule
A misses blocks for a long time, no MSIG & $\beta$ does not change\\
matching update whose new map leaves only A ON & close the old version; NextLive may select A again\\
matching update whose new map has A OFF and B ON & close the old version and hand off to B\\
matching update whose new map sets every region OFF & keep $G$ and the target index in a PENDING tail\\
later update targeting that pending version with B ON & create another version and hand off to B\\\bottomrule
\end{tabularx}
\end{center}

\subsection{Repeated Force and an unassigned boundary}
\label{sec:repeated-force}
A valid MSIG can close a newly assigned version while its target is
still Assigned, or after it starts but before any DATA. Forces need no DATA
in between; their thresholds need only be $T_{\rm update}$ apart. Each
new instruction still needs a matching
target and a valid advance announcement.

If closing session $(7,20)$, which has accepted output, finds no eligible
region, replay stores target index 21 and anchor $G$ in version 8. It
opens no regional engine. A later valid MSIG can move the version to 9
while keeping target index 21 and the same $G$. If that update's complete
map sets B ON, B enters $(9,21,0,0)$, not $(9,22,0,0)$: a session that
never started does not count as a completed turn. If no successful update
ever supplies an eligible map, the wait can last forever.

\subsection{Strict cutoff, equal timestamps and concurrent preparation}
\label{sec:cutoff-example}
An MSIG targets version 7 at threshold 300. B may start regular session
$(7,21)$ from A's terminal receipt before that threshold. Meanwhile its
baseline state follows $G$, so it already holds the state at any Force
anchor. The baseline grants no entry.

At the first finalized block with timestamp at least 300, the prelude
checks the still-current version. A match closes every session of
version 7 and keeps the exact preceding $G$. Old DATA in that block's
payload cannot revoke the control. A program created for that
Force starts from the state at that exact anchor, never from state
computed for another one.

If Espresso timestamps jump across the threshold, the first crossing
block is still the activation block. A later block with the same
timestamp does not repeat an already processed instruction. A distinct
valid instruction can execute at a later position if its own target
matches. The model excludes a timestamp jump that would make two valid
updates activate in the same block (\cref{ass:single-activation}); such a
history is outside the model of this document.

\subsection{Conditions for progress}
\label{sec:progress}
Being Ready only means that a region can start. Global progress further
needs regional consensus, a FINAL receipt, publication and global
acceptance. The following conditions are separate from prefix safety.
\begin{itemize}
\item Regional service keeps Part~I's fair leader schedules, eventual
exclusion of exposed equivocators, actual entry into correct turns, and
uninterrupted commit opportunities. These are conditions inside a region,
not fairness among regions.
\item Espresso keeps finalizing retrievable history with growing
timestamps. Complete results reach the Publish Layer, and the declared
publication and finalization services hold.
\item The required authority, exact anchor state and acceptance evidence
are supplied. A selected live region satisfies \cref{ass:ready-service}.
The next required DATA is published in a form and position that pass the
current-branch acceptance checks.
\item If the current region cannot serve, recovery additionally depends
on the multisignature service actually submitting a suitable atomic
update whose complete map keeps or enables an eligible working region,
and eventually allowing
enough uninterrupted execution for regional commitment and delivery.
\end{itemize}
The last item is a liveness condition on the control service, not a
condition for accepting an MSIG. The protocol specifies no detection or
response bound, the only interval it enforces between Forces is the
Espresso-time gap $T_{\rm update}$, and $T_{\rm L1}$ starts only at actual
submission. So an honest service can still fail to
give a useful execution opportunity unless its policy provides one.

A wholly offline region is not an operational region for these service
claims. Every accepted certificate still satisfies the fixed-snapshot,
lifetime key-compromise and historical-validation premises. If Espresso fails, global confirmation stops; switching the
finality service is outside this protocol.

\subsection{What is needed for global progress}
\label{sec:global-progress}
Under the Part~I conditions above, a session that keeps running can form
a complete result. The bounded preparation assumption alone does not
provide that uninterrupted regional opportunity.

If a complete result is available at time $s$, the declared services give
submission by $s+T_{\rm pub}$ and Espresso visibility by
$s+T_{\rm pub}+T_E$. This is transport of the result. To advance $G$, the
result must supply the next required DATA under its fixed boundary, while
its producing branch version is still open at the canonical acceptance
position, and its references and predecessors must pass the checks of
\cref{sec:exports}. Because the rule ignores outer hints, reads the
cumulative $\mathcal O_q$ and appends one DATA at a time, wrong hints or
a later pending successor do not erase an otherwise valid next DATA. They never permit a gap or an append from a
closed version.

With these regional, transport, acceptance and control-service
conditions, replay of the suitable result advances $G$. The service
bounds give no wall-clock processing deadline for every arbitrary
observer. No unconditional global-growth theorem follows from repeated
Ready events, published raw records or committee honesty alone.

\paragraph{Offline recovery requires control.}
When the current region cannot continue, progress needs the committee to
submit a suitable atomic update whose complete map disables a failed
candidate when necessary and keeps or enables a live eligible successor.
If no such update is
submitted, nothing else forces a switch. If all regions are OFF, the
boundary can stay unassigned forever. If $T_{\rm update}$ is shorter than
a useful interval, MSIGs can keep closing every one, and no regional
output need ever become global. These executions
are allowed unless the governance service provides more liveness
conditions. The committee can thus stall the chain, but under the
safety assumptions it cannot use a Force to roll $G$ back.

\paragraph{Failure and safety are separate.}
An offline region does not meet the operational premises of the
preparation and regional-service claims. A violated lifetime Byzantine
weight bound is a different failure, of the certificate model. Neither
case lets the multisignature committee roll back $G$. No result here
promises fair service among regions, automatic ON, one DATA between
Forces, or unconditional global progress.

\Needspace{14\baselineskip}
\begingroup\small
\begin{thebibliography}{9}
\bibitem{part1}
Sync HotStuff RR, Part I: Core Protocol. Design document, 13 October 2026.
Companion source: \href{run:sync_hotstuff_rr_part1_core.tex}{Part I core protocol}.
\bibitem{part3}
Sync HotStuff RR, Part III: Staking and Authority. Version 3, 19 October 2026.
Companion source: \href{run:sync_hotstuff_rr_part3_staking_v3.tex}{Part III staking and authority}.
\bibitem{part4}
Sync HotStuff RR, Part IV: Force Handler and Publish Layer.
Version 0.1, 20 October 2026.
Companion source: \href{run:sync_hotstuff_rr_part4_force_v0_1.tex}{Part IV Force handler and Publish Layer}.
\bibitem{sync}
I. Abraham, D. Malkhi, K. Nayak, L. Ren, and M. Yin.
\emph{Sync HotStuff: Simple and Practical Synchronous State Machine Replication}.
IEEE Symposium on Security and Privacy, 2020, pp. 106--118.
\href{https://conferences.computer.org/sp/pdfs/sp/2020/349700a106.pdf}{Published paper}.
\bibitem{certification}
MegaETH Specification, PR~8, \emph{Block Certification}.
Commit \texttt{aedc2826663b}.
\href{https://github.com/megaeth-labs/specification/tree/aedc2826663b134ee25d26bea17e70c6d23d321c/specs/block-certification}{Pinned specification}.
Its Espresso and certifier verification interfaces are a reference;
the rotation and explicit-control rules in this document are new.
\end{thebibliography}
\endgroup
\appendix
\crefalias{section}{appendix}
\section{Security and lifecycle guarantees}
\label{app:security}
This appendix collects the public FINAL-kernel, startup, replay,
program-isolation, explicit-control and preparation propositions.
Their assumptions and qualifications are stated with each result.

\subsection{Soundness of the public FINAL kernel}
\label{app:final-kernel-proof}
\begin{proposition}[Soundness of the public FINAL kernel]
\label{prop:final-kernel-soundness}
Let $C$ be the authentic incoming context of session $(\beta,r)$, and
suppose $\tagg{RawFINAL}_C(D,\Sigma)$ holds. Then some correct replica of
that session committed exactly this DATA $D$ through its own commit
timer, and before doing so it fully validated $D$ and the inputs $D$
depends on. Hence $D$ is a validly admitted regional input under its
anchor, including the $\rho$ and $\xi$ conditions, whatever witness a
publisher happens to attach. The same holds for every same-session DATA
ancestor of $D$ on its hash-linked path.
\end{proposition}

We restate \cref{prop:final-kernel-soundness} in the precise form the
proof uses.

\paragraph{Precise statement.}
\textit{Assume the fixed-snapshot fault bound, the accepted-lifetime
key-compromise bound, the cryptographic assumptions, sound validation
interfaces, and the regional full-validation and FINAL-generation rules.
If $C$ is the authentic incoming context of $(\beta,r)$ and
$\tagg{RawFINAL}_C(D,\Sigma)$ holds, then a correct replica committed
this exact DATA $D$ at a timer event, directly or as an ancestor. Before
that event it fully checked its own-session proposal and required input
ancestry, using established historical kernel conclusions only for
inherited anchors. So a genuine valid historical witness exists for $D$.
Consequently the exact DATA has valid historical regional admission and
native input validity under its authenticated anchor, including the
regional $\rho$ and $\xi$ validity conditions. Exact same-session DATA
ancestors on that hash-linked path inherit these commitment and
input-validity conclusions.}

\begin{proof}
By the fixed-snapshot fault bound, a set of distinct signers with weight
greater than $F_i$ includes a correct signer. The lifetime compromise
condition and signature unforgeability make its signature an actual
correct FINAL action on the stated DATA hash and context. In particular,
the signature binds the exact initial anchor in $C$, so evidence from a
different incoming boundary for the same pair cannot pass this check.
This step uses the raw threshold, not the definition of a valid regional
receipt.

A correct replica sends FINAL only after Commit's timer event commits the
named DATA, either as the timed DATA itself or in the ancestry of the
timed DATA/DRAIN block. Vote started that timer only after all
current-session dependencies were present and the proposal and required
ancestry passed full validation. In particular the replica checked its
session's historical authority and admission, parent justifications,
quota and positions, the L1 and Espresso reference rules, and native
input validity. The sound interfaces apply to these actual regional
checks. An inherited anchor from another producer may rely on that
producer's historical kernel conclusion, justified by the induction
below; this does not replace the replica's exact native state or its
own-session validation. These facts give a genuine valid historical
witness for the committed DATA. It need not be the witness that a later
publisher attaches.

A commit triggered by a receipt, including Completion, emits no new
FINAL. So it cannot supply the first correct signature without the
earlier fully validated timer event identified above. Dependencies on
earlier accepted anchors are justified by induction over execution
events from the trusted initial checkpoint. The FINAL signatures used by
an inherited anchor exist before the new Vote event, so their own timer
origins are earlier cases of the induction. The same argument covers an
unaccepted REGULAR anchor without assuming its global acceptance. No step
assumes the validity of the kernel acceptance being proved. Hash binding
identifies the supplied DATA and exact parent-linked bodies with those in
the correct signer's validated history. The timer's ancestor-commit rule
gives the same-session coverage claim.
\end{proof}

\subsection{Synchronous start}
\begin{proposition}[Synchronous start]
\label{prop:sync-start}
Let $e$ be the first correct initial Enter for $(\beta,r)$ and write
$\Delta=\Delta_{i_{\beta,r}}$. Suppose the handoff remains applicable,
the correct members stay operational, and validation, state adoption
and Enter are covered by the complete-message processing bound.
Then every correct selected member enters by $e+\Delta$.
\end{proposition}
\begin{proof}
The first correct entrant sends a complete ReadyBundle at $e$.
Every correct recipient validates it and completes guarded Enter by
$e+\Delta$, unless it entered earlier.
\end{proof}

\subsection{Conditional safety and lifecycle invariants}
\label{sec:safety-argument}
The two propositions below cover public replay and the local programs.
They are safety statements only; progress is discussed in
\cref{sec:liveness}. They use the model and verification interfaces of
\cref{sec:model}. In plain words: everyone who reads the same complete
history accepts the same chain; DATA in $G$ is never removed by a Force;
and no program enters twice or mixes state with another session.

\paragraph{Regional agreement before Force.}
For a branch that has not been forced, the fixed-handoff initialization
has the same base case as Part~I's genesis: a common anchor, one
initial authority and entry skew at most $\Delta$. The first DATA uses
$Q^0_{\beta,r}$; later DATA and all recovery generations keep Part~I's
parent and lock rules. This is conditional on the ReadyBundle contract
of \cref{prop:sync-start} and the regional validation and fault premises.

Completion does not silently remove a voter from a surviving commit
window. Suppose a correct timer started at $s$ expires at $s+2\Delta$.
If another correct member completed by $s+\Delta$, its complete terminal
receipt would reach this member by expiry and cancel that timer, with
message processing first at equality. Quit has the same cancellation
argument. Initial entry finishes before the earliest possible terminal
receipt, since entry takes at most $\Delta$ and a direct timer needs
$2\Delta$. Thus the ordinary vote-window argument applies to every
surviving direct timer before Force: matching correct votes are cast by
$s+\Delta$, and their protecting QC arrives by $s+2\Delta$. Recovery
locks then preserve that commitment as in Part~I. Receipt-based commits
trace to earlier direct timer commitments and add no new FINAL
signatures, so induction over those events also covers Completion.
Successive regular handoffs inherit the resulting common terminal DATA.

This argument stops at Force activation. It is not used to claim
agreement among lagging nodes on a discarded suffix. A correct node
cannot verify that Force before its finalized activation block exists;
the earlier public input sets therefore contain only evidence generated
before such forced termination. The public cutoff, proved next, excludes
old-version output from the activation block onward.

\begin{proposition}[Common interpretation and preservation of global DATA]
\label{prop:global-prefix}
Let two correct observers start from the same authenticated replay
checkpoint and apply \cref{alg:global-rule} under the stated history,
announcement and verification premises. Their global DATA sequences are
prefix-comparable: one is a prefix of the other. Replaying to the same canonical subposition with all applicable source
evidence and L1 controls verified yields the same output
$\mathrm{Out}=(G,\mathsf{Tail})$ of \cref{sec:exports}. At a completed Espresso prefix
this gives the corresponding completed-prefix result as a special case. It does not require identical
future-control queues or source-reading checkpoints. Every accepted DATA
block keeps its content, position, producing identity and bound
session record forever. Once a version is closed, no
later EXPORT adds new DATA from that version to $G$.
\end{proposition}
\begin{proof}
Induct over the common finalized Espresso sequence and, within a block,
over the steps of \cref{alg:global-rule}.
The initial checkpoints agree. Complete applicable L1 controls,
immutable authorities, Espresso headers and namespace bytes authenticate
the same next input. Advance announcement supplies controls before their
crossing; missing actual evidence pauses at the same operation. Atomic
empty-or-singleton activation, the target test, target indices and fixed-map
NextLive selection are deterministic. They give the same prelude and
immutable snapshots. Future controls not yet due change neither result.

Both observers read the same object set $\mathcal O_q$ from the same
authenticated namespace content. A private cache changes retrieval work,
not $\mathcal O_q$. The matching conditions of Step~2 and its hash tie-break are
functions of $\mathcal O_q$ and the current output, so both observers
append the same block. By \cref{prop:final-kernel-soundness},
\tagg{CheckFINAL} establishes intrinsic input validity from a genuine
correct signer's regional commit. The independent boundary, ancestry,
source-order and cutoff checks establish current append permission.
Common input and the specified tie-break identify the same next DATA;
this step does not assume regional agreement during forced termination.
Optional admission or
native sidecars cannot veto that kernel. Executable state, native output
and DA correspondence remain subject to their separate sound interfaces. Each DATA transaction appends the
same next block and accounting, and persists the same subposition.
A terminal transaction constructs the same boundary and applies the same
fixed-map selector. It either selects the same eligible successor, returns
the same pending tail, or pauses at the same missing evidence. Missing evidence
cannot undo the already committed transaction or alter its timestamp.
These rules give agreement by induction over all committed subpositions.
Two observers that have reached different subpositions hold ordered
states on this append-only DATA sequence, including when one is paused
inside a source record. Extra future-control queues and private validation
caches are not part of the compared replay result.

Inspect preservation at each step. Control ingestion removes no DATA. A
successful update keeps all of $G$ at the start of the block and changes
only the branch version, its map and the tail for future appends. A
pending boundary also preserves $G$. Each accepted DATA transaction only extends it. Retained DATA keeps its
original producing pair and hash. Closed versions are never reopened,
and every later DATA candidate checks the current version. Thus
neither an old regular successor nor a late receipt can append a closed
suffix. For a Force at $J_{q_f}$, earlier canonical appends are retained,
and Step~1 rejects the old version before processing $M_{q_f}$.
A late observer replaying an earlier block cannot use evidence from a
later $\mathcal O_q$. Thus late old votes or receipts cannot alter
either the retained prefix or its historical replay. Local termination
does not mutate $G$.
\end{proof}

\begin{proposition}[Fixed entry and isolation of local programs]
\label{prop:program-isolation}
Under the regional-validity and source-verification premises, each
correct replica installs at most one handoff for a pair $(\beta,r)$.
Its initial Enter occurs at most once, uses exactly that installed
anchor and state, and never uses an obsolete task callback as authority.
Actions for one pair do not alter another pair's locks,
votes, timers or entry time. Baseline execution creates none of
these BFT actions.
\end{proposition}
\begin{proof}
Proceed from the initial checkpoint through successive handoffs.
For a regular handoff, the predecessor's fixed boundary and quota fix
its terminal height. A verified terminal receipt and its exact ancestry
fix the proposed successor's anchor and signed Espresso reference.
Before Force, regional agreement makes that terminal common among
correct members. After Force, discarded histories need not agree, but
the local record still cannot be replaced. The immutable $S_\beta$ and
target authority
fix the region, while the terminal reference fixes the next opening
floor. For Force, common replay fixes the old and new versions, source
context, exact retained $G$ and newly installed map; NextLive under that
map fixes the target. Repeated Forces without assignment
change $\beta$ but keep the stored predecessor index. They cannot
reuse an already opened pair. Handoff verification rejects conflicting
fixed fields instead of replacing an existing record.

Baseline state follows $G$ over time but is not a program. Only Create
creates a program, and only after verification installs the actual
handoff.
Material callbacks keep the requested pair and anchor identity;
the gate checks the exact fixed state and separate current-activation predicate.
The serialized Enter call tests and sets the durable entry flag
before dispatching signatures. Repetition, restart and late callbacks
therefore cannot repeat Enter or reset its entry time. Each core
message is authenticated and routed by the pair it signs; local tasks
use the same pair and anchor context. An Assigned program does not run
the core.
Version-wide closed marks and Terminated records prevent a late event
from reopening an old program. Force-close applies to every old pair,
including one first named by a later packet. These checks establish
the stated isolation by induction
over the local event sequence.
\end{proof}

This is local authorization relative to a node's processed public
prefix. It does not claim that all nodes learn a finalized Force at the
same real time. A lagging node can briefly run an old program, but
\cref{prop:global-prefix} still excludes its closed-version output.
The protocol makes no regional agreement claim for that discarded
suffix during forced termination. The public FINAL kernel still proves
a correct signer's historical validation and commit, with the initial
anchor hash bound in its signed context. It does not grant append
permission to the closed version.
\subsection{Explicit control and conditional preparation}
\label{sec:liveness}
The guarantees in this subsection do not require every region to receive
service.

\begin{proposition}[No unannounced cross-region Force]
\label{prop:explicit-force}
Starting from an authenticated checkpoint, every later increment of
$\beta$ has a unique effective MSIG cause in the verified L1 control
history. Missing DATA, a local health report or a Ready event does not
increment $\beta$; there is no separate status-update transition.
\end{proposition}
\begin{proof}
Step~1 of \cref{alg:global-rule}, detailed in Step~3 of
the Force handler in Part~IV~\cite{part4}, is the only transition that increments $\beta$.
It is reached only for the unique update activated at the block
whose target matches the branch version recorded at its start. Control identities
and persisted subpositions prevent a second application. Other
transitions select under an already fixed map, append DATA or update
local programs; none changes $\beta$ or installs a map. Induction over committed
transitions proves the claim.
\end{proof}

\begin{proposition}[Preparation under a continuing assignment]
\label{prop:ready-bound}
Let a public assignment open in block $J_o$, with finalization time
$t_0$. If the selected region and its state service satisfy
\cref{ass:ready-service}, the fixed handoff stays applicable, and no
Force supersedes it during preparation, then correct target members can
complete Enter by $t_0+T_{\rm ready}$. Entry uses the retained
anchor of that assignment and happens at most once per member and pair.
Entry does not require any DATA production.
\end{proposition}
\begin{proof}
A member that already entered satisfies the claim. For a member that has
not entered, the service assumption supplies every material condition of
\cref{def:ready} by the stated time, and continuing applicability
supplies the current-entry check. \Cref{alg:rotation,alg:bootstrap} then
perform the guarded entry in the same complete event.
\Cref{prop:program-isolation} gives anchor binding and at-most-once
entry. The gate does not look at future DATA output. A valid Force before
completion would remove the continuing-assignment premise; it would not
violate any waiting rule.
\end{proof}

\section{A reference control policy}
\label{app:policy}
This appendix sketches one policy that an honest committee can run on
top of the interface of \cref{sec:runtime-controls}. The protocol does
not depend on it, and certifiers do not check it. Other conditions or
thresholds keep every guarantee stated in this document.

\paragraph{What the controller observes.}
The controller acts on a verified replay position $q$. There it reads
$\mathrm{Out}_q$, the current immutable map $S_\beta$ and the
control log, including its pending updates. It also knows the complete
regional results it has received, and operational reports. When the
$\mathsf{Tail}_q=\tagg{NEXT}(\nu)$, write
$(\beta,r,i,a)=(\beta(\nu),r(\nu),i(\nu),a(G_q,\nu))$
for the export context. On the global clock of \cref{ass:network},
the controller records $t_0$ when it first verifies this export
assignment and $t_G$ when it accepts its latest DATA. Let $t$ be its
current time. These local observation times need not equal a block's
real finalization time; $t_0$ is no earlier than the public opening used
by \cref{ass:ready-service}.

\paragraph{Conditions.}
The policy uses the following predicates, with policy parameters
$\tau_{\rm boot}$, $\tau_{\rm stall}$ and $k_{\rm off}\ge1$.
\begin{itemize}
\item $\tagg{Pub}$: the controller knows a complete regional result for
DATA height $h(\operatorname{last}(G))+1$ of the export session that is
not yet in $\mathcal O_q$.
\item $\tagg{Boot}$: $\mathsf{Tail}_q=\tagg{NEXT}(\nu)$, $a=0$, $\tagg{Pub}$ fails,
and $t-t_0>\tau_{\rm boot}$.
\item $\tagg{Stall}$: $\mathsf{Tail}_q=\tagg{NEXT}(\nu)$, $a>0$, $\tagg{Pub}$ fails,
and $t-t_G>\tau_{\rm stall}$.
\item $\tagg{Fault}(i)$: the controller holds evidence that region $i$
cannot be trusted or cannot serve. The evidence is \emph{checkable}, for
example an admissible signed leader equivocation, or
\emph{attested}, for example an operator or security report. Divergence
in a suffix discarded by Force is not by itself evidence that an honest
member violated this protocol.
\item $\tagg{Maint}(i,t_m)$: region $i$ is to stop serving at time $t_m$.
\item $\tagg{Restore}(i)$: region $i$ is OFF, and the controller holds
evidence that it can serve again.
\item $\tagg{Unassigned}$: $\mathsf{Tail}_q=\tagg{PENDING}(\omega)$.
\end{itemize}
$\tagg{PendU}$ holds when the controller has an issued update that has
not yet reached its activation block. $\mathit{strikes}(i)$ counts the
Forces issued under \tagg{BOOT} or \tagg{STALL} whose export region was
$i$, since $i$ last added DATA to $G$.

\paragraph{Updates issued.}
The reference policy is \defterm{single-flight}: while $\tagg{PendU}$
holds, it issues no second update. It accumulates desired operational
changes off chain. After the outstanding entry is canonically processed,
the controller rereads the resulting $\beta$ and $S_\beta$, reevaluates
all conditions, and, if action is still needed, submits one update that
targets that $\beta$ and carries one complete new map. This avoids a
queue in which the first update changes $\beta$ and makes the rest stale.

When no update is in flight, start from $S^\circ=S_\beta$, apply every
condition that currently holds, and batch the result into one complete
$S_{\rm new}$. \Cref{tab:control-conditions} gives the intended map
change. Every issued update also forces a branch change; there is no
status-only row.

\begin{table}[htbp]\centering\small
\caption{Reference policy: conditions and the update issued.}
\label{tab:control-conditions}
\begin{tabularx}{\linewidth}{@{}>{\raggedright\arraybackslash}p{.15\linewidth}>{\raggedright\arraybackslash}p{.2\linewidth}YY@{}}\toprule
Code & Condition & Complete $S_{\rm new}$ & Replay effect\\\midrule
\tagg{FAULT} & $\tagg{Fault}(i)$ & $S^\circ$ with $i$ OFF & Close
$\beta$, retain $G$, and select under the new map\\
\tagg{MAINT} & $\tagg{Maint}(i,t_m)$ & $S^\circ$ with $i$ OFF, with
activation chosen for $t_m$ & Close $\beta$ at activation and select
under the new map\\
\tagg{BOOT} & $\tagg{Boot}$, not $\tagg{PendU}$ & $S^\circ$, with export region $i$
OFF if $\mathit{strikes}(i)+1\ge k_{\rm off}$ & Close
$\beta$; target index $r$; selection after $i$\\
\tagg{STALL} & $\tagg{Stall}$, not $\tagg{PendU}$ & As for \tagg{BOOT} &
Close $\beta$; target index $r+1$; selection after $i$\\
\tagg{PUB} & $\tagg{Pub}$ & No update; ask a backup sender to publish
(Part~IV~\cite{part4}) & None\\
\tagg{RESTORE} & $\tagg{Restore}(i)$ & $S^\circ$ with $i$ ON & Close
$\beta$ and make $i$ eligible on the new branch\\
\tagg{UNASSIGNED} & $\tagg{Unassigned}$ & $S^\circ$ with a chosen region
$j$ ON & Close the pending version; keep its target index and try $j$ on
the new branch\\
none & None of the above, including an L1 or Espresso outage & No update
& None\\\bottomrule
\end{tabularx}
\end{table}

Under this interface even \tagg{RESTORE} and maintenance of an idle
region close the current branch, because every map change is atomic with
a Force. The policy should therefore batch such changes rather than use
the interface for frequent health telemetry. A \tagg{FAULT} with
checkable evidence calls for an update at the earliest admissible
activation, subject to the single-flight rule. $\tagg{Boot}$ and
$\tagg{Stall}$ require that $\tagg{Pub}$ fails on purpose: when the next
DATA exists but has not been published, a Force would discard committed
output and gain nothing.

\paragraph{Issuance rules.}
\begin{enumerate}[label=(R\arabic*)]
\item \textbf{Complete map and target.} Every update names the branch
version the controller observed and carries a complete map derived from
that same replay state. A pending update is never rewritten. If it turns
out stale, its entire effect is lost and the controller reevaluates the
conditions against the new current version.
\item \textbf{Single flight.} Normally issue no update while one is
unactivated. After it executes or becomes stale, resubmit any still
desired intent with a fresh sequence, target and complete map.
\item \textbf{Timing.} At actual submission time $s$, choose the
earliest admissible threshold: above $s+T_{\rm L1}+\epsilon_E$
(\cref{sec:assumptions}) and at least $T_{\rm update}$ after the last
issued threshold. The update takes effect in whatever state replay has
reached by then. If submission slips, sign a fresh update.
\item \textbf{Short horizon.} Do not schedule an update more than a
horizon $T_{\rm horizon}$ ahead; keep planned changes off-chain until they are near.
A far-future in-flight update delays any later intervention under the
single-flight policy.
\item \textbf{Keep a candidate.} Do not issue a map that leaves no
eligible region, unless the code is \tagg{FAULT} with checkable
evidence.
\item \textbf{No anchor.} The controller never tries to pick the
retained DATA; replay keeps the actual $G$ at the start of
$J_{q(\theta_a)}$.
\end{enumerate}

\paragraph{Choosing the thresholds.}
\tagg{BOOT} and \tagg{STALL} should never fire against a region that
serves as the model allows. Let $T_{\rm first}$ bound the time from
$t_0+T_{\rm ready}$ until a complete result for the session's first DATA
is available, and $T_{\rm next}$ the time from one complete result of
the session until the next. These bounds come from the regional service
conditions of Part~I; this document does not derive them. Choose
\[
 \tau_{\rm boot}\ge T_{\rm ready}+T_{\rm first}+T_{\rm pub}+T_E,
 \qquad
 \tau_{\rm stall}\ge T_{\rm next}+T_{\rm pub}+T_E.
\]

\begin{proposition}[The policy does not interrupt a serving region]
\label{prop:policy-noninterference}
Let the thresholds satisfy the bounds above, and let the controller keep
up with Espresso history as in \cref{ass:services}(iii), has the
required acceptance evidence, and processes each eligible result when
that evidence arrives. Consider an assigned export session that no
Force closes, whose region satisfies
\cref{ass:ready-service}, whose first complete result is available by
$t_0+T_{\rm ready}+T_{\rm first}$, whose later results each follow the
previous one within $T_{\rm next}$, and whose results pass the
acceptance checks of \cref{sec:exports}. Then $\tagg{Boot}$ and
$\tagg{Stall}$ never hold for this session.
\end{proposition}
\begin{proof}
The first result is available by $t_0+T_{\rm ready}+T_{\rm first}$. By
\cref{ass:services}(i) and (iii), it is in a finalized Espresso block,
and verified by the controller, by
$t_0+T_{\rm ready}+T_{\rm first}+T_{\rm pub}+T_E\le
t_0+\tau_{\rm boot}$. Replay accepts it in that block, because
it extends $G$, passes the acceptance checks and its branch version is
open. So once $t-t_0>\tau_{\rm boot}$, the controller already
sees $a>0$, and $\tagg{Boot}$ fails. For $\tagg{Stall}$, let the last
accepted DATA be appended at $t_G$. Its result was available at some
time $s\le t_G$. The next result is available by $s+T_{\rm next}$ and is
accepted and verified by $t_G+T_{\rm next}+T_{\rm pub}+T_E\le
t_G+\tau_{\rm stall}$. So $t_G$ moves forward before $\tagg{Stall}$ can
hold. When the session's last DATA is accepted, the tail no longer names
this session, and neither predicate applies to it.
\end{proof}

The proposition covers only \tagg{BOOT} and \tagg{STALL}. A \tagg{FAULT}
or \tagg{MAINT} decision, or a threshold set too low, can still
interrupt a serving region. While $G$ is fixed, each \tagg{BOOT} or
\tagg{STALL} Force moves the selection to the next eligible region, and
a serving region is never interrupted by them; this is the step that
recovery needs in \cref{sec:progress}. A full recovery theorem also needs
the authority of each new target index to keep that region eligible; we
do not prove it here.

\section{Node storage and source relay}
\label{app:storage}
\label{sec:source-step}
\label{sec:source-retention}
A node must keep the exact bytes of old public objects. The object set
$\mathcal O_q$ of \cref{alg:global-rule} only grows, so the rule may
append a DATA whose body or FINAL signatures arrived many blocks
earlier; a handoff or recovery may also refer to old objects. Keeping
only a hash is not enough. This appendix states what each regional
member stores, how it starts storing, and how it relays new source data
to its region. In short: keep every byte that may be needed later, store
it before moving on, and relay new source data to the region at once.

\paragraph{The public store.}
Each correct regional member keeps a durable \tagg{PublicStore} shared
by its source reader, replay logic, Session Manager and engines. It stores content, not only cursors and indexes. The store
has the following distinct components.
\begin{table}[htbp]\centering\small
\begin{tabularx}{\linewidth}{@{}>{\raggedright\arraybackslash}p{.24\linewidth}Y@{}}\toprule
Component & Stored content\\\midrule
Raw Espresso source & Every acquired complete relevant namespace,
including rejected or malformed records; headers, finality/linkage
proofs, authenticated absence, and canonical source order.\\
Object set $\mathcal O_q$ & The block bodies and FINAL signatures
published through $J_q$, with their bytes, indexed by hash and session,
and the source block that carried each. An index without the bytes is
not a stored object.\\
L1 source objects & Every obtained finalized configuration, recovery record, control,
header and registry/completeness proof, with its authentication inputs;
pending controls as well as applied controls.\\
Regional witnesses & Full validation bundles actually used by correct
regional actions, committed DATA and receipt dependencies. Possession
does not put a private object into $\mathcal O_q$.\\
Prepared anchors & Exact executable continuation objects and their
validated provenance, separately keyed by anchor, configuration,
region and native format.\\
Replay state & Verified checkpoint,
$\mathcal X=(\mathrm{Out},S_\beta,\kappa)$ of \cref{alg:global-rule}, the replay subposition,
source-availability cursors, and typed pending fetches.
The source-availability cursor is distinct from the replay cursor.\\\bottomrule
\end{tabularx}
\end{table}

At each source-availability event, store the entire authenticated raw
namespace and its evidence before advancing the availability cursor,
releasing the input buffer, or making a relay depend on it. An object
that \cref{alg:global-rule} has not used, or never uses, is stored all
the same. Likewise store each acquired L1
object before advancing its authority-fetch cursor. Regional validation
stores the witness it used, so a later required complete send does not
reconstruct that witness from an archive. Correct FINAL and complete
receipt sends obey the full-bundle rule of \cref{sec:admission}; a
reference is used only when its receiver can resolve it within the
regional delay bound.

A new branch version, session closure, advancement of $G$, checkpoint
compaction and program exit do not delete these objects. Cancellation
removes a consumer, not the store's retention responsibility. Transfer
outstanding sends and ownership durably before releasing resources.
Restart restores the exact store and replay subposition. If required
bytes are lost, the node must hydrate them again before the timing
claims for prepared nodes apply to it again.
These are protocol storage rules, not an additional storage service.

\paragraph{Finite initial hydration and pruning.}
\label{sec:source-hydration}
Choose a fixed authenticated source cut $q_0$ when beginning this
retention discipline. Hydrate the raw source interval still legally
referenceable after that cut and all exact objects designated by the
starting checkpoint. Because $\mathcal O_q$ is cumulative, this
includes relevant raw namespaces back to the authorized source origin,
including objects that were never appended. Storing only accepted
ancestry is insufficient. Also prepare the fixed anchor's provenance and the
authority and state dependencies of the assigned handoff.
Store every newly available namespace beyond $q_0$ while hydration
runs. Neither a new branch version nor a later source height restarts the fixed
old task or changes its target to a moving cut.

The interval and each admitted payload are finite. Eventual historical
retrieval gives a finite completion time for that fixed set; the anchor
and finite L1 source prefix have their separate eventual-availability
premises. Taking a finite maximum imposes no bound on arbitrary archive catch-up.
The selected-region preparation service is stated separately in
\cref{ass:ready-service}. It does not include all future
L1 objects that an arbitrary sidecar might name.

A replay checkpoint may compact computation without deleting
referenceable bytes. Deduplication and equivalent lossless storage are
allowed while exact content and source-occurrence mappings remain
available to the modeled handler. A hash or an only-eventually-readable
archive is not a replacement for those bytes. Pruning requires a
protocol proof that no future supported reference, checkpoint,
anchor, pending validation or publication duty can need the object.
Advancing $G$ or terminating its producing engine does not prove this.
\Cref{alg:global-rule} itself never appends an object of a closed
version, or a DATA at or below the height of $\operatorname{last}(G)$
that is not in $G$; other duties may still need such objects, so this
alone does not justify pruning them.
The present rule supplies no finite-storage bound.

\paragraph{What the rule reads from the store.}
\label{sec:record-completeness}
\Cref{alg:global-rule} reads only $\mathcal O_q$ and the L1 inputs of
\cref{sec:logical-carrier}. $\mathcal O_q$ consists of the block bodies
and FINAL signatures carried by the authenticated namespace content of
$J_1,\ldots,J_q$, together with the objects that a protocol-authorized
checkpoint establishes. A private cache, a newer private checkpoint, a
hash or name without its object, or an unresolved sender assertion adds
nothing to $\mathcal O_q$. L1 objects keep their own authentication
interface and need not first occur in Espresso.

For a candidate $B$ in Step~2, the rule needs exactly the body of $B$,
the bodies of a covering path to a root $D$, the FINAL signatures on
$D$, and the descriptor $\nu$ in the tail $\tagg{NEXT}(\nu)$ that
acceptance will bind to
$B$. A publisher-selected PROPOSE,
recovery $\tagg{AUTH}_{\beta,r,g}$ path, L1 reference proof sidecar or unbound
native claim is not an input; its optional verification never delays the
rule. The correct signer's full regional checks justify intrinsic input
validity by \cref{prop:final-kernel-soundness}; they do not turn a
private object or a sidecar assertion into a public input. Check typed
identities, actual hashes, the signed domain, the current pair and the
fixed snapshot before verifying FINAL signatures and applying
\tagg{CheckFINAL}. The outer fields of an EXPORT are not checked.

An absent required object or an invalid signature only means that the
candidate does not qualify yet. Other candidates are unaffected, and
since $\mathcal O_q$ only grows, the candidate can qualify in a later
block once the missing object is published. An old object stays in
$\mathcal O_q$ even if nothing was appended when it arrived. It never
grants append permission to a closed version. In contrast, a missing
complete namespace, applicable control prefix or initial authority for
the exact target pair of an atomic update or terminal successor selection
pauses the rule at that operation. It is never zero
weight or an empty candidate set. A cache miss for an object already in
$\mathcal O_q$ requires retrieval and resumption at the same
subposition. For a continuously online member that follows these
retention rules, \cref{prop:retained-public-bytes} rules out such a miss
of Espresso bytes. Optional regional evidence is
stored and routed to its own engine; a complete genuine Quit keeps its
evidence-first handling.

Step~2 appends, persists and exposes the terminal DATA before its
terminal branch applies NextLive for the successor. A wait for successor
authority therefore cannot withhold that DATA. Packet preparation cannot
invert this order.

\paragraph{Immediate regional source relay.}
A \tagg{SOURCESTEP} is an authenticated message to an installed fixed
regional snapshot. It carries its chain/protocol domain and sender,
the source height and full header, complete relevant namespace or
absence, finality/linkage evidence and canonical record order. It
identifies the maintained verified predecessor checkpoint, supplies
applicable complete L1 controls and the available fixed-boundary
inputs, and supplies required exact public bytes not already resolvable
by the recipient under the complete-message contract. The
recipient never trusts a sender-selected new checkpoint.

On acquiring a complete raw source input, each correct member stores
and sends it in that availability event, before waiting for its own
public replay result under \cref{alg:global-rule}. Add its regional objects to the store
without first fetching optional witness targets. This relay
need not know which region replay will select. Native preparation does
not delay it. Missing mandatory boundary evidence remains an explicit
wait and is supplied in a later complete dependency event. An initial
raw relay is not falsely described as a complete boundary operation
when that evidence is absent. Equal packets are deduplicated by source
position and commitment without suppressing the first correct relay.

A sender includes full required bytes unless recipient resolution is
already covered by the regional full-bundle contract. Its own cache hit
alone does not justify omitting bytes at a recipient. Acquisition of an
additional exact dependency immediately resumes its saved operation;
there is no independent polling tick or replay job.

\paragraph{Receive, commit and activate.}
The receiver authenticates and stores the input, then processes the
contiguous dependency-complete prefix from its maintained checkpoint.
Out-of-order source inputs wait for their predecessors. For each ready
block it executes \cref{alg:global-rule}. Each step records its change to the output and
its subposition, and exposes appended DATA before any later wait. A
duplicate cannot repeat a completed step, count or boundary. Local
receive times do not replace canonical timestamps.

The same complete event delivers the verified transition to the Session
Manager. When the exact anchor continuation and remaining full Ready
inputs are already prepared, the Session Manager checks the fixed handoff and
the Ready condition and performs initial Enter once. It
attaches the already executable object; it does not import a database
or reconstruct native state here. No old lock, vote or entry time is
inherited. A Force-created version still cannot export in its triggering
block.

On verified closure, persist the version-wide closed mark and apply
Force-close in \cref{alg:rotation}. Every BFT action, including a timer
callback, tests that mark before it runs. Unknown old pairs are blocked
by the same mark, so no membership discovery or close relay is needed
to stop them. Required bytes and already owed sends remain in the shared
store; cancelling the program does not cancel their retention. Physical
cleanup neither delays the cutoff nor blocks entry into a valid new
session.

\begin{proposition}[Stored public source bytes]
\label{prop:retained-public-bytes}
After fixed-cut hydration, a continuously online history-following
member that executes the retention rules keeps every legally
referenceable acquired Espresso object and its authentication evidence.
If $J_q$ first becomes final at real time $t_E(q)$, the required
Espresso bytes through $q$ are available to that member by $t_E(q)+T_E$, without a
new archive read for a previously acquired object.
\end{proposition}
\begin{proof}
Hydration establishes the claim through the fixed cut. At each later
availability event, insertion of the entire raw namespace precedes the
availability-cursor advance and any use by \cref{alg:global-rule}.
Earlier objects of $\mathcal O_q$ are already stored; new objects reside
in this namespace. An absent
object is not fabricated. No retry, checkpoint or program transition
removes referenceable content. Induction over source positions preserves
the claim. Every predecessor through $q$ became final no later than $t_E(q)$,
so the stated visibility premise supplies its new source bytes and
Espresso evidence by $t_E(q)+T_E$. Older hydrated bytes remain present.
This is a statement about Espresso bytes, not a deadline for new L1 source records.
\end{proof}

\paragraph{Prepared regional processing.}
\label{sec:source-step-timing}
When source bytes, authority and exact executable continuation are
resident, the complete receive handler validates them, performs the
canonical public transition and checks the guarded initial entry.
No historical fetch or native import is silently charged to the
regional communication bound.

The complete-message contract bounds delivery and processing of its
declared inputs. Initial preparation has the separate scope and origin
of \cref{ass:ready-service}. Returning-node catch-up outside that scope
remains eventual. The source-availability cursor, canonical replay
cursor and local Ready event can differ; none supplies authority to
execute an unannounced Force.

\end{document}
